\documentclass[anonymous,notoc]{gmcm2026}
\usepackage{placeins,seqsplit,calc,needspace,float}
\providecommand{\tightlist}{\setlength{\itemsep}{0pt}\setlength{\parskip}{0pt}}
\providecommand{\pandocbounded}[1]{#1}
\setlength{\emergencystretch}{2em}
\makeatletter\setlength{\@fptop}{0pt}\makeatother
\renewcommand{\topfraction}{.90}
\renewcommand{\bottomfraction}{.85}
\renewcommand{\textfraction}{.08}
\renewcommand{\floatpagefraction}{.84}
\setcounter{topnumber}{3}
\setcounter{bottomnumber}{2}
\setcounter{totalnumber}{4}
\setlength{\textfloatsep}{12pt plus 2pt minus 2pt}
\setlength{\intextsep}{10pt plus 2pt minus 2pt}
\setlength{\floatsep}{10pt plus 2pt minus 2pt}
\widowpenalty=10000
\clubpenalty=10000
\displaywidowpenalty=10000
\renewcommand{\qedsymbol}{\ensuremath{\scriptstyle\square}}
\gmcmsetup{title={多核神经网络处理器的计算图切分与调度优化},year=2026,edition={第二十三届},problem=A}
\begin{document}
\maketitle
\begin{abstract}
神经网络处理器（neural processing unit，NPU）包含多个计算核心（NPU core，图中简称 Core）。待处理的计算任务包含一组运算、数据搬运以及它们的先后关系。将操作和数据分别表示为节点，用有向边表示操作读取或生成哪些数据，就得到本题的计算图。本文把图中的计算操作划分为若干组，连同相关数据和依赖构成原图的子图；各核心实际执行的是这些子图所对应的计算任务。因此，需要共同决定操作怎样分组、各组任务交给哪个核心，以及任务按什么顺序执行。

问题一规定，每个子图对应一个独立下发的执行任务（Task），任务之间传递的数据必须经过外部主存。算法寻找互不需要对方结果的两组运算。在满足后文给出的任务高度、接收核心任务顺序及外部依赖条件时，将其中一组移交另一核心，并把需要读取两组结果的后续运算单独安排。任务顺序若相互冲突，可能导致各任务都在等待尚未完成的其他任务。因此，本文先证明所给调整条件能够避免这种循环等待，再比较调整前后的总体任务执行时间（Makespan）。

问题二把分配到同一个核心的各子图所对应的操作合并到一个Task中，允许后续操作直接使用该核心缓存中保留的数据。算法同时考虑生成、读取同一数据的操作应放在哪些核心，以及计算单元何时空闲：前者用于计算复制字节数，后者用于安排计算位置与次序。问题三进一步增加各核心共享的只读缓存。算法调整计算顺序，以改变结果需要保留的时间和后续读取的时刻；对于中间结果只被唯一父操作读取的树状计算，证明在逐棵子树完整处理的限定模型下，所给排序使内部结果的最大存储量最小。

三个固定算法版本均完成100张官方计算图、每图1～5个核心的评价，每题得到500个合法结果。问题一在2～5个核心上的平均加速比为1.9516、2.7514、3.4786、4.0553；问题二为2.3167、3.2356、3.9711、4.5498。问题三在5个核心上的平均基线加速比为4.7576；保持调度方案不变，启用只读缓存相对关闭缓存的逐图平均加速比为1.0829，逐图字节命中率均值为31.36\%。问题一5个核心时的平均求解时间为3.522秒，问题二全部测试的平均求解时间为4.029秒，均包含生成和比较方案所需的时间。

实验中也出现复制量减少但完成时间变长的方案。因此，本文先用依赖和存储关系提出调整，再按官方程序对完整方案的模拟结果选择输出；复制字节数和局部存储量只承担其定义范围内的判断，不代替最终完成时间。
\keywords{多核调度；计算图切分；数据搬运；缓存复用；总体任务执行时间}
\end{abstract}
\section{问题重述与分析}

\subsection{研究对象与优化目标}

本题面向神经网络处理器（neural processing unit，NPU）的多核调度优化问题。NPU 包含多个计算核心（NPU core，下文正文统称为``核心''，图表中简称为 Core）。NPU 代表整颗处理器芯片系统，而核心是指其内部并行的独立计算与控制单元。每个核心配备专用的矩阵计算单元 Cube、向量计算单元 Vector、核内核外数据搬运单元，以及该核心独占访问的私有片上缓存------局部存储器 L1 与统一缓冲区（Unified Buffer，简称 UB）。各核心之间无法直接相互读写对方的私有存储空间，所有核心通过总线共享核外主存储器，题面简称为 DDR（源自双倍数据率动态随机存取内存，double data rate dynamic random-access memory）。DDR 的总读写带宽、片上私有缓存容量以及各类硬件操作之间的同步等待周期均由硬件系统固定。

交由 NPU 执行的计算任务，在底层由一系列具有前后依赖关系的算子运算及其读写的数据组成。一个算子通常需要等待其前驱算子输出所需的数据后方可启动。为了形式化地刻画这些运算任务及其时序与数据流动约束，题面将待执行的神经网络模型表示为\textbf{计算图（computational graph）}。算法的目标是在给定的硬件约束下，共同决定计算任务的操作划分、核心分配以及执行次序{[}1{]}。

以两次相互独立的矩阵乘计算为例，将它们分别分配至两个核心可以实现时间上的并发执行；然而，若两项计算均需从核外 DDR 载入大规模输入张量，它们就会在总带宽有限的 DDR 访问通道上产生争用。此外，若后续运算必须等待这两次矩阵乘全部结束方能汇总，系统还必须计入数据传输延迟与跨核心同步等待。由此可见，单纯平衡各核心的计算量只是调度优化的基础前提，完成时间的长短更取决于数据搬运开销、共享带宽争用以及依赖等待链的综合作用。

本题的首要评价指标为总体任务执行时间（Makespan），即在官方离散事件模拟评估程序下，从所有操作的最早启动时刻到全部操作均完成的最晚结束时刻所经历的硬件时钟周期数（cycle）。与方案质量指标相对应的是算法本身的求解效率，即算法程序在通用计算设备上从读取输入图到持久化输出合法方案所耗费的物理墙钟时间（Wall-clock Time），本文统称为求解时间（以秒为单位）。前者衡量调度方案在模拟环境下的执行性能，后者反映算法的实用工程效率。本文报告的执行时刻与周期数均源自官方评测系统的模拟推进，不作为物理硬件测量的绝对性能指标。

\subsection{计算图的数据表示与数据依赖}

题面将神经网络抽象为二部有向无环图（bipartite directed acyclic graph，bipartite DAG）形式的计算图。图中的节点集合被严格划分为两类互不相交的顶点：算子节点（Op node，表示一次具体的数学计算或数据搬运）与数据张量节点（Tensor node，表示输入数据、计算中间结果或最终输出数据）。

有向边仅连接不同类别的节点，表达数据的生成与消费逻辑：若算子 \(u\) 运算生成数据张量 \(t\)，则计算图中存在有向边 \(u \to t\)；若算子 \(v\) 需要读取张量 \(t\) 作为计算输入，则存在有向边 \(t \to v\)。图中不存在算子到算子或张量到张量的直接连边。张量节点包含 \texttt{size} 属性，记录该数据块所占用的物理存储大小（以字节 Byte 为单位）。

根据计算图的有向拓扑，操作之间存在两类依赖关系：若算子 \(u\) 生成张量 \(t\)，且算子 \(v\) 直接读取该张量 \(t\)（即二部图中仅通过单一中间张量节点连接 \(u \to t \to v\)，且不经过任何其他计算操作），则称操作 \(v\) \textbf{直接数据依赖}于操作 \(u\)；此时，后继操作 \(v\) 的开始时刻不得早于前驱操作 \(u\) 的完全完成时刻。若从操作 \(u\) 到操作 \(v\) 的有向路径中途经过了一个或多个其他中间操作（即形成了 \(u \to \cdots \to w \to \cdots \to v\) 的传递链条），则称操作 \(v\) \textbf{间接数据依赖}于操作 \(u\)。只要两个操作之间存在直接或间接数据依赖，后继操作 \(v\) 的开始时刻就必须滞后于前驱操作 \(u\) 的结束时刻。若两个计算操作仅共同读取同一份外部输入张量，它们之间并不存在先后生产关系，因而在逻辑上不构成相互数据依赖。计算图的有向无环特性确保了计算流单向推进，不存在自相矛盾的时序闭环。

在后续算法设计与拓扑分析中，为了清晰展示计算任务的先后偏序关系，我们常将二部图向仅包含算子节点的\textbf{算子依赖图}进行投影：对二部图中每一个张量节点 \(t\)，若算子 \(u\) 是生成 \(t\) 的操作（\(u \to t\)），且操作集合 \(\{v \mid t \to v\}\) 读取 \(t\)，则在算子依赖图中由 \(u\) 向该集合中的每一个读取算子 \(v\) 引出一条有向依赖边 \(u \to v\)。算子依赖图是用于简化依赖推导的辅助表示，底层的存储占用、搬运开销以及缓存判定仍基于带有张量大小属性的原始输入图展开。

根据官方输入输出契约，题面预设了两类核外数据搬运操作：\texttt{COPY\_IN} 负责将指定张量从核外 DDR 搬运至对应核心的私有缓存（L1 或 UB）；\texttt{COPY\_OUT} 负责将计算核心内部私有缓存中的数据写回至 DDR。算子的唯一标识编号（Op ID）与张量的标识编号（Tensor ID）属于两个独立的命名空间，分别唯一确定底层的调度单元与数据对象。

\subsection{执行资源与计算和数据搬运的时间重叠}

在每个核心内部，物理硬件被组织为若干专用的硬件执行管道（Pipe），题面将其运行机制统称为\textbf{硬件流水线（pipeline）}。在本文所针对的 NPU 体系结构语境下，硬件流水线（pipeline）特指核心内部服务于特定指令类型的微架构执行资源及其严格保序的处理机制。根据题面硬件建模，每个核心包含四条相互平行的 Pipe： 1. \texttt{PIPE\_M}：驱动矩阵计算单元 Cube，专门负责矩阵乘累加等密集矩阵算子； 2. \texttt{PIPE\_V}：驱动向量计算单元 Vector，专门负责向量算术、非线性激活等向量算子； 3. \texttt{PIPE\_MTE2}：驱动核内数据搬运单元，负责将数据从核外 DDR 搬入私有缓存 L1/UB，或在私有缓存内将数据从 UB 转移至 L1； 4. \texttt{PIPE\_MTE3}：驱动核内数据搬运单元，负责将私有缓存 L1/UB 中的数据写回至 DDR，或在私有缓存内从 L1 转移至 UB。

官方离散事件评测程序对流水线推进作出了明确的时序约束：在单个核心内部，同一条 Pipe 在任意时刻至多只能执行一个操作；不同 Pipe 之间只要满足数据依赖和片上存储容量条件，则允许在时间轴上并发推进。

例如，当核心内的向量流水线（\texttt{PIPE\_V}）正在对已位于私有缓存中的数据执行运算时，搬运流水线（\texttt{PIPE\_MTE2}）可以同时从核外 DDR 载入后续算子所需的其他输入张量。两项操作如果在时间轴上拥有一段共同的执行区间，便在逻辑上达成了\textbf{计算与数据搬运的时间重叠（overlap of computation and data transfer）}。在模拟模型中，实现时间重叠能够有效减少处理器的空转等待。能否顺利实现时间重叠，取决于输入数据是否已传输到位、对应执行 Pipe 是否处于空闲状态、所属任务是否获准激活，以及私有缓存 L1/UB 中是否有足够的可用存储容量以容纳输出张量。若因等待前驱结果或缓存不足导致 Pipe 空闲停顿，系统时间线上就会产生调度气泡（Bubble）。

\subsection{计算操作的分组、分配与调度方案表示}

面向多核心体系架构，算法需要协同确定三项核心决策： 1. \textbf{操作划分}：将计算图中全部非核外搬运的待调度算子划分为若干互不重叠的集合，每个集合连同其引用的局部张量及依赖边构成原计算图的一个\textbf{子图（subgraph）}； 2. \textbf{核心分配}：将划分出的各个子图分别映射指派至特定的 NPU 核心负责处理； 3. \textbf{时序编排}：在同一个核心内部，确定该核心所承载的各个子图之间的执行先后次序。

根据官方提交契约（\texttt{contest\_io.py}），提交方案必须且仅包含 \texttt{node\_to\_subgraph} 和 \texttt{core\_schedules} 两个字段。其中 \texttt{node\_to\_subgraph} 映射字典负责记录原计算图中操作类型不属于 \texttt{COPY\_IN} 与 \texttt{COPY\_OUT} 的全部算子所属的子图编号（核内搬运或其它非核外 COPY 算子仍须纳入划分）；\texttt{core\_schedules} 列表为每个核心分别给出该核心所分配子图编号的执行顺序序列。方案必须保证上述全部待调度算子被无遗漏、无重复地覆盖，且每个子图编号在所有核心的调度列表中恰好出现一次。

在评测系统的执行语义中，\textbf{执行任务（Task）}是分配给特定核心执行的计算单元。子图是算法在计算图拓扑层面划分的操作集合，而 Task 则是评测模型运行时的调度与隔离边界，核心实际执行的是子图所对应的计算任务。赛题所提出的三个问题，其本质差异正是由子图到 Task 的组织关系以及片上数据驻留规则的改变所界定的：

{\def\LTcaptype{table} % do not increment counter
\begin{table}[!htbp]\centering\small
\begin{tabular}{@{}
  >{\raggedright\arraybackslash}p{(\linewidth - 6\tabcolsep) * \real{0.2500}}
  >{\raggedright\arraybackslash}p{(\linewidth - 6\tabcolsep) * \real{0.2500}}
  >{\raggedright\arraybackslash}p{(\linewidth - 6\tabcolsep) * \real{0.2500}}
  >{\raggedright\arraybackslash}p{(\linewidth - 6\tabcolsep) * \real{0.2500}}@{}}
\toprule\noalign{}
\begin{minipage}[b]{\linewidth}\raggedright
机制与规则维度
\end{minipage} & \begin{minipage}[b]{\linewidth}\raggedright
问题一
\end{minipage} & \begin{minipage}[b]{\linewidth}\raggedright
问题二
\end{minipage} & \begin{minipage}[b]{\linewidth}\raggedright
问题三
\end{minipage} \\
\midrule\noalign{}


Task 的组织构成 & 每个子图独立构成一个 Task & 同一核心所分配的全部子图合并为一个统一的 Task & 与问题二相同 \\
同核跨子图数据保留 & 严禁保留，必须经由 DDR 写回与重新载入 & 允许在私有缓存 L1/UB 中跨子图直接驻留与复用 & 与问题二相同 \\
同一核心内切换延迟 & 连续 Task 间必须等待 100 cycles & 核心仅执行单一合并 Task，无内部 Task 切换延迟 & 与问题二相同 \\
跨核心通信同步延迟 & 跨核前驱 Task 完成后需等待 1000 cycles & 对应源 \texttt{COPY\_OUT} 完成后需等待 500 cycles & 与问题二相同，即使命中只读 Cache 也需满足该时序约束 \\
共享二级缓存（L2） & 无，核心仅能访问私有 L1/UB 与共享 DDR & 无，核心仅能访问私有 L1/UB 与共享 DDR & 引入核外共享只读缓存（Level-2 Cache，简称 L2） \\
\bottomrule\end{tabular}
\end{table}
}

由上述规则演进可知： - 问题一将每个子图视为隔离的 Task，算法面临增加多核并行度与增加 DDR 搬运及 Task 切换边界之间的直接冲突； - 问题二通过在同一核心内将子图合并为大 Task，解除了同核跨子图的数据强制落盘限制，允许中间结果在私有片上缓存中直接驻留复用，但中间张量在缓存中停留越久，对私有存储空间的占用压力就越大； - 问题三在多核合并执行的基础上引入了全局只读 Cache，多个核心读取相同张量时，若数据已驻留于 Cache 则产生\textbf{缓存命中（cache hit）}，可直接从高速 Cache 通道读取而无需重复向 DDR 发起请求，从而有效分流共享带宽；若不在 Cache 中则发生\textbf{缓存未命中（cache miss）}，需从 DDR 搬入并在搬入完成后按先进先出（FIFO）规则写入 Cache；单张量大小超过 Cache 总容量时则不予缓存。

图 \ref{fig:progression} 针对相同的代表性计算片段，直观对比了三类问题在 Task 构成、数据流动路径以及多核心交互方式上的结构演变。

\begin{figure}[!htbp]
\centering
\includegraphics[width=\linewidth,height=.84\textheight,keepaspectratio]{figures/fig-progression-v3.png}
\caption{三问条件的差别。问题一为图中不同子图所对应的操作分别建立Task，数据经DDR中转；问题二允许同一个核心保留数据；问题三另为符合条件的COPY\_IN提供共享只读Cache。下部示意数据能够装入Cache时的读取过程；所有请求仍须满足原依赖及规定等待。}
\label{fig:progression}
\end{figure}

\subsection{多核调度中的核心挑战与求解思路}

综合题意与硬件执行特征，在通用神经网络处理器上实现高效的多核调度存在以下三项核心难点：

第一，\textbf{并行计算收益与跨核通信开销之间的非线性权衡}。将计算图拆解为更多子图并分散至更多核心，理论上能够提升计算单元的宏观并发度；然而，跨核心的数据交换必须插入额外的 \texttt{COPY\_OUT} 与 \texttt{COPY\_IN} 搬运操作。所有核心的 DDR 搬运操作均受制于总带宽仅为 60 字节/周期的共享内存通道，密集的跨核通信极易引发严重的内存带宽争用与同步停顿。如果因通信导致的额外 DDR 搬运与等待所增加的执行周期，超过了并行分配所节省的计算周期，多核方案的总完成时间反而会恶化。

第二，\textbf{片上私有缓存容量限制与数据驻留的动态约束}。在问题二和问题三中，将具有数据依赖的操作集中至同一核心能够避免外部内存搬运；然而，核心内部 L1 和 UB 的私有存储容量是严格受限的。核内调度在编译阶段若检测到局部缓存容量不足，会预先插入将数据换出至核外 DDR、后续再换入的溢出搬运操作（spill）。这不仅会增加额外的 DDR 搬运开销，若调度依赖安排不当还会导致方案无法推进。

第三，\textbf{局部核心排序与全局数据依赖交织下的死锁风险}。通过反例，我们可以说明为什么仅仅保证原始计算图无环并不能确保多核调度方案合法：假设核心 0 上的子图任务 \(S_1\) 需等待核心 1 上的子图任务 \(S_2\) 提供数据，而核心 1 内部指定的调度次序又要求先执行任务 \(S_3\) 再执行任务 \(S_2\)；若任务 \(S_3\) 恰好又间接依赖于任务 \(S_1\) 的运算结果，两个核心便会陷入``\(S_1\) 等待 \(S_2\)、\(S_2\) 等待 \(S_3\)、\(S_3\) 等待 \(S_1\)''的循环依赖状态，导致评测系统发生死锁。

针对上述挑战，本文遵循题面关于避免暴力迭代搜索、单用例高效求解的指导原则，设计了结构化的求解流程： 1. \textbf{基于图结构的特征分析与候选构造}：分析计算图的连通分支、关键路径与张量复用关系，针对独立并行分支与数据密集型计算链，直接构造少量具有结构特征的完整候选方案； 2. \textbf{调度顺序的无环性判定}：将核心内子图执行次序作为时序边显式引入子图依赖图中，证明所给次序不产生有向环，确保方案满足依赖先后关系（需要说明的是，顺序无环是合法执行的必要条件，但仍需结合存储容量与完整评测进一步检验）； 3. \textbf{结合通信代价模型的理论下界快速剪枝}：基于题面给定的硬件流水线资源与通信带宽约束，计算理想并发条件下的时间下界与跨核通信开销，辅助快速筛除劣质候选； 4. \textbf{离散事件模拟选优}：在受限的评估预算内，调用评测程序对生成的完整候选方案进行离散事件模拟，从中选取 Makespan 最小的方案作为输出。


\section{模型假设与符号说明}

\subsection{硬件参数与执行约束}

本文的建模与算法设计严格建立在题面给定的体系结构参数与官方离散事件评测程序的基础上。计算图中各算子的计算语义、操作类型及原始数据依赖保持不变，算法不得通过删减计算或修改运算结果来缩短完成时间。全部存储容量、总线带宽及同步等待常数均采用官方固定配置{[}1{]}。

根据硬件配置参数，系统的主要物理约束如下： 1. \textbf{核心数量}：可用核心数 \(K \in \{1, 2, 3, 4, 5\}\)，系统为各核心编号为 \(0, 1, \dots, K-1\)； 2. \textbf{片上私有缓存}：每个核心配备独立的局部存储器 L1 与统一缓冲区 UB，其存储容量分别为 \(524{,}288\text{ B}\)（\(512\text{ KiB}\)）与 \(131{,}072\text{ B}\)（\(128\text{ KiB}\)）。两级缓存各自独立管理，未分配的 L1 容量无法与 UB 共享； 3. \textbf{共享主存与缓存带宽}：所有核心访问核外主存（DDR）的读写操作共享统一的全局带宽池，总带宽固定为 \(60\text{ B/cycle}\)。在问题三中，引入容量为 \(1{,}048{,}576\text{ B}\)（\(1\text{ MiB}\)）的共享只读二级缓存（Level-2 Cache，简称 L2），命中 Cache 的读取操作使用专用的共享 Cache 读取带宽通道，总带宽为 \(250\text{ B/cycle}\)，该带宽与 DDR 带宽相互独立，不占用 DDR 带宽池； 4. \textbf{同步与切换延迟}：在问题一中，同一核心连续执行的相邻 Task 之间须等待 \(100\text{ cycles}\) 的切换时间；若一个 Task 依赖位于其他核心的前驱 Task，其启动时刻还须满足晚于前驱 Task 完成时刻至少 \(1000\text{ cycles}\) 的跨核同步延迟。在问题二和问题三中，由于同核子图合并为单一 Task，核心内部不存在 Task 切换延迟；跨核心依赖仅作用于具体数据搬运操作，目标核心的 \texttt{COPY\_IN} 需在源核心对应的 \texttt{COPY\_OUT} 完成后等待 \(500\text{ cycles}\) 方可参与就绪判断。这些常数由评测程序在对应的事件发生时检查，不直接等同于图上的固定边权重。

评测程序采用离散事件模拟推进执行。对于流水线操作队列，算子严格按核内静态调度生成的偏序队列以先进先出（first in, first out，FIFO）方式推进就绪与发射。对于问题三中的只读 Cache，其缓存行维护独立的 FIFO 淘汰队列：当 Cache 剩余容量不足以容纳新载入的张量时，系统依次淘汰最早写入 Cache 的条目；后续核心读取 Cache 发生命中时，只读访问不会更新该条目在 FIFO 队列中的淘汰顺序。单张量物理大小超过 Cache 总容量（\(1\text{ MiB}\)）的数据不进入 Cache。

若核内调度算法在静态编译时评估到私有片上缓存（L1 或 UB）容量不足以容纳计算过程中的所有活跃张量，评测程序会在该核心的指令流中预先插入溢出搬运操作（spill），即通过 \texttt{PIPE\_MTE3} 将暂不使用的中间数据经由 DDR 写出，后续需要时再通过 \texttt{PIPE\_MTE2} 重新从 DDR 读入。溢出搬运会产生额外的 DDR 访问流量并占用共享带宽，调度方案应尽量平抑片上张量驻留峰值以减少溢出开销。

\subsection{原始张量、数据副本与逻辑标识}

在计算图的生命周期中，同一份数学数据在不同调度阶段可能对应不同的图节点与存储实体，我们明确区分以下三类对象： 1. \textbf{原图张量节点}：指原始输入计算图 JSON 文件中给定的张量实体，具有唯一的整型 \texttt{id} 和固定的数据大小 \texttt{size}； 2. \textbf{调度生成的数据副本}：在多核心切分与核内调度阶段，当数据发生跨核心传输或因容量受限触发溢出时，评测程序会为该数据在不同存储介质（如源核心私有缓存、核外 DDR、目标核心私有缓存）中派生出新的物理数据实例。这些副本虽然承载相同的内容，但在物理上分别占用各自存储介质的独立空间； 3. \textbf{逻辑数据标识（\texttt{logical\_tid}）}：评测系统在处理问题三的只读 Cache 查询时，为所有源于同一张量派生出的物理副本赋予统一的逻辑标识 \texttt{logical\_tid}。当核心发起 \texttt{COPY\_IN} 搬运时，评测程序依据请求数据的 \texttt{logical\_tid} 查询 Cache 队列。若命中则直接从 Cache 读取，若未命中则在 DDR 读取完成后以该 \texttt{logical\_tid} 登记至 Cache 中。

此外，方案合法性必须满足三重判定条件： - \textbf{操作覆盖合法}：输入图中除 \texttt{COPY\_IN} 与 \texttt{COPY\_OUT} 之外的全部算子，均被无重复、无遗漏地指派至唯一的子图，且每个子图在 \texttt{core\_schedules} 中恰好出现一次； - \textbf{依赖拓扑无环}：子图间的数据依赖与各核心指定的子图执行顺序合并构成的扩展依赖图不存在有向环； - \textbf{存储与执行可行}：评测程序对方案能够成功完成核内调度并推进离散事件模拟，不因容量冲突或死锁导致模拟异常终止。

只有评测程序完整运行成功并输出日志的方案，其导出的 Makespan 才是有效指标。若方案因违反依赖或容量超限被评测程序拒绝，或者执行发生超时，均视为无效解，不将其填为零参与均值统计。

\subsection{主要符号}

本文使用的主要数学符号定义如下表所示。

{\def\LTcaptype{table} % do not increment counter
\begin{longtable}[]{@{}
  >{\raggedright\arraybackslash}p{(\linewidth - 4\tabcolsep) * \real{0.3333}}
  >{\raggedright\arraybackslash}p{(\linewidth - 4\tabcolsep) * \real{0.3333}}
  >{\raggedright\arraybackslash}p{(\linewidth - 4\tabcolsep) * \real{0.3333}}@{}}
\toprule\noalign{}
\begin{minipage}[b]{\linewidth}\raggedright
符号
\end{minipage} & \begin{minipage}[b]{\linewidth}\raggedright
定义说明
\end{minipage} & \begin{minipage}[b]{\linewidth}\raggedright
取值范围或量纲
\end{minipage} \\
\midrule\noalign{}
\endhead
\bottomrule\noalign{}
\endlastfoot
\(G=(O\cup T, E)\) & 原始计算图，其中 \(O\) 为操作节点集合，\(T\) 为张量节点集合，\(E\) 为有向依赖边集合 & 题面输入给定 \\
\(V\) & 待调度的计算算子集合（即 \(O\) 中操作类型不为 \texttt{COPY\_IN} 和 \texttt{COPY\_OUT} 的算子） & \(V \subset O\) \\
\(b_t, \ell_t\) & 数据张量 \(t\) 的物理大小（字节）及其预设存储位置（DDR、L1 或 UB） & \(b_t \in \mathbb{N}^+ \text{ (B)},\ \ell_t \in \{\text{DDR}, \text{L1}, \text{UB}\}\) \\
\(d_v, p_v\) & 算子 \(v\) 的硬件执行周期数及其所绑定的硬件执行管道 Pipe & \(d_v \ge 1 \text{ (cycle)},\ p_v \in \{\text{PIPE\_M}, \text{PIPE\_V}, \dots\}\) \\
\(K\) & 系统可用的计算核心总数 & \(K \in \{1, 2, 3, 4, 5\}\) \\
\(\pi(v)\) & 算子划分映射，表示算子 \(v\) 所归属的子图编号 & \(\pi: V \to \mathbb{N}\) \\
\(a(s)\) & 子图分配映射，表示子图 \(s\) 所指派执行的核心编号 & \(a: \mathcal{S} \to \{0, 1, \dots, K-1\}\) \\
\(\sigma_k\) & 核心 \(k\) 所承担的子图编号的有序调度列表 & 子图编号的全排列序列 \\
\(P=(\pi, a, \sigma)\) & 完整的多核切分与调度方案（对应提交 JSON 文件的两个字段） & 参赛提交规范 \\
\(q\) & 赛题问题编号 & \(q \in \{1, 2, 3\}\) \\
\(M_q(P)\) & 方案 \(P\) 在问题 \(q\) 的评测环境下运行所得的总体任务执行时间（Makespan） & cycle \\
\(D_q(P)\) & 方案 \(P\) 在评测程序下产生的总额外 DDR 数据搬运量 & B（字节） \\
\(B_i\) & 第 \(i\) 张计算图在官方基准单核心调度算法下的基线 Makespan & cycle \\
\(H_{i,K}\) & 第 \(i\) 张计算图在 \(K\) 个核心下运行问题三方案所得的 Cache 字节命中率 & \([0, 1]\) \\
\(L, U\) & 在相同图输入、核心数与硬件配置下，对所有合法方案均成立的理论下界与已有合法方案的时间上界 & cycle \\
\(L(P)\) & 仅在指定候选方案 \(P\) 的特定结构约束下成立的时间下界（不作为全局下界） & cycle \\
\(U/L\) & 当前最佳方案的 Makespan 与理论下界之比，用于评估当前方案逼近理论极限的程度 & 无量纲比值 \\
\(T_{\text{solve}}\) & 求解算法从读取输入到输出合法方案文件的端到端墙钟运行时间 & s（秒） \\
\end{longtable}
}

在官方输入文件中，算子的 \texttt{cycles} 字段可能为 0。根据官方评测程序逻辑，实际调度推进时各算子至少占用 1 个执行时钟周期。因此在建立数学分析模型与计算下界时，算子执行周期均按 \(d_v = \max(1, \text{cycles}_v)\) 计算。

\subsection{评价指标及统计方法}

\subsubsection{1. 方案质量指标与平均加速比}

多核心调度方案的质量首要由 Makespan \(M_q(P)\) 衡量。为了评估多核并行化相对于单核执行的性能提升程度，根据题面第 5 页脚注 2 的明确规定，多核心调度方案的平均加速比（Speedup）采用逐用例比值的算术平均方法进行计算：

\[\bar{S}_{q,K} = \frac{1}{100} \sum_{i=1}^{100} \frac{B_i}{M_q(P_{i,K})},\]

其中 \(B_i\) 为官方评测系统针对第 \(i\) 张图在单核心条件下的基线 Makespan，\(P_{i,K}\) 为算法针对第 \(i\) 张图在 \(K\) 个核心下生成的调度方案。必须说明的是，平均加速比是逐图加速比值的算术平均值，而非全部用例单核时间总和除以多核时间总和。题面官方基线 \(B_i\) 由通用的单核启发式调度算法给出，并不保证是单核条件下的数学全局最优解，因此在部分用例中可能出现加速比略大于核心数 \(K\) 的情况。

对于问题三，题面要求评估引入只读 Cache 带来的相对性能增益。评测时保持问题三所生成的调度方案 \(P\) 完全不变，分别在启用 Cache 和关闭 Cache（即回退至问题二纯 DDR 模式）两种配置下运行评测程序，分别记录其 Makespan 为 \(M_3(P)\) 和 \(M_2(P)\)。单用例的 Cache 加速比定义为 \(M_2(P) / M_3(P)\)，多用例汇总时同样报告这 100 对结果比值的算术平均值。

\subsubsection{2. 额外数据搬运量与缓存命中率}

次要方案质量指标包括数据搬运量与缓存命中率： - \textbf{总额外 DDR 数据搬运量 \(D_q(P)\)}：官方评测程序统计调度后执行的所有 \texttt{COPY\_IN} 和 \texttt{COPY\_OUT} 搬运字节总数，减去原始计算图固有输入输出张量所需的基准搬运字节数。在问题三中，官方评测系统保持相同的统计口径，并不在 Cache 命中时将该统计字段直接扣减，因此该字段反映的是切分引入的总通信需求，而非物理 DDR 总线的实际净流量； - \textbf{Cache 字节命中率 \(H\)}：定义为命中 Cache 的数据总字节数占所有发起 Cache 查询的总数据字节数的比例，即 \[H = \frac{\text{hit\_bytes}}{\text{hit\_bytes} + \text{miss\_bytes}}.\] 当用例中不存在任何可查询的 \texttt{COPY\_IN} 操作时，官方程序将分母为零的情形约定为 \(H = 0\)。该指标严格按字节量统计，而非按访问请求的发生次数统计。

\subsubsection{3. 求解效率与原题时限解读}

算法的求解效率通过端到端求解时间 \(T_{\text{solve}}\) 评估，度量从程序启动、解析输入图 JSON，到完成拓扑分析、方案构造、内部筛选及输出方案文件全过程的实际物理墙钟耗时（Wall-clock Time）。

原题第 5 页脚注 1 指出：``此处`高效地'代表着算法应该避免使用暴力迭代搜索机制，以单个用例能在 5～10 分钟内给出结果为佳。''我们对该项要求的理解为： - 该项表述属于针对求解算法工程效率与设计理念的指导性建议，强调算法应当利用计算图的固有结构与硬件特性进行针对性构造，避免盲目的大规模蒙特卡洛枚举或无剪枝的暴力迭代； - 5～10 分钟是算法在常规计算硬件上的用例级运行耗时建议，并非超过 600 秒即被判定无效的硬性淘汰线，亦非要求算法必须耗满 5 分钟方可输出。本队算法在保证方案高质量的前提下，追求极高的求解速度，单用例求解耗时通常控制在数秒至数十秒之间。

\subsubsection{4. 评测程序的角色与调用账目}

在算法研发与评测交付中，涉及三套评测与打分程序： - \textbf{E0（官方离散事件评测程序）}：官方提供的未修改基准评测脚本。本文汇报的所有主实验成绩、最终 Makespan 以及附录逐用例结果，均以 E0 的运行输出为最终准绳； - \textbf{E1（批量精确评估工具）}：队内依据官方规则实现的轻量级独立评估实现，用于问题一中在受限调用预算内（单用例至多 10 次）对少量候选方案进行离散事件精确评分与选优； - \textbf{E2（代理代价评估工具）}：队内面向快速探索开发的代理代价模型。在最终冻结的主求解器中，问题二与问题三均已实现基于图结构与数学理论的直接构造，求解过程中无需依赖在线迭代评分（问题二直接利用超图最小割模型与约束构造方案并由 E0 确认，问题三基于归约树排序理论生成方案并由至多 3 次 E0 调用验证）。


\section{计算图切分与调度模型}

\subsection{计算图划分、任务分配与执行顺序的数学表示}

设\(V\)为原图中全部需要划分的计算操作。函数\(\pi\)为每个操作指定子图编号，所有非空子图组成集合\(\mathcal S\)；函数\(a\)为每个子图对应的任务指定核心；列表\(\sigma_k\)给出核心 \(k\)上各子图对应任务的先后顺序。三者合在一起构成方案\(P=(\pi,a,\sigma)\)。仅有节点分组还不是完整方案，因为同一分组采用不同核心分配或不同执行顺序，可能得到不同结果。

如果子图\(s\)中的某个操作需要子图\(t\)中的操作产生的数据，就在子图之间画\(t\to s\)。全部这类边组成\(E_{\mathcal S}\)。此外，若同一核心的列表中\(s\)紧接在\(t\)之后，还需画\(t\to s\)表示执行顺序；这些边组成\(E_\sigma\)。因此，用于检查先后关系的图为

\[D^+(P)=(\mathcal S,E_{\mathcal S}\cup E_\sigma).\]

本文将同时包含\(E_{\mathcal S}\)和\(E_\sigma\)的\(D^+(P)\)称为\textbf{加入同一核心顺序后的子图依赖图}。如果这些关系形成有向环，就无法为相应任务安排满足所有关系的先后顺序，方案不能采用。

通过顺序检查后，官方程序按问题编号\(q\)组织Task，补充边界搬运，再插入必要的缓存换出、换入和同步关系，最后模拟执行。将这两个过程分别记为\(\mathcal C_q\)和\(\mathcal R_q\)，则

\[P\xrightarrow{\mathcal C_q}\mathcal E_q\xrightarrow{\mathcal R_q}(M_q,D_q,H_q).\]

\(\mathcal E_q\)表示组织后的操作、依赖和资源约束，\(H_q\)只在问题三报告。相同大小的两次DDR搬运可能因同时存在的其他请求不同而耗时不同，因此不能在这个阶段为每个Task预先指定一个与其他Task无关的固定时长。

令\(\mathcal F_q(G,K)\)为官方程序能够成功执行的完整方案集合。主要优化问题为

\[\min_{P\in\mathcal F_q(G,K)}M_q(P).\]

搬运量\(D_q(P)\)和问题三的Cache命中率另行报告。问题一、二比较方案时优先比较完成时间，时间相同再比较规定的搬运指标；问题三只有在完成时间严格下降时才替换原方案。上述规则是本文的算法选择规则，不是题目给出的加权评分公式。

\subsection{三个问题的统一求解过程}

三问的共同困难在于：便于计算的局部指标不能完整描述机器执行，但对每一种可能划分都进行完整模拟又过于昂贵。本文按以下顺序求解。

第一，从原图识别有明确含义的关系。例如，哪些计算互不依赖，哪些操作读取同一张量，哪些中间结果只供唯一后续操作使用。第二，针对这些关系生成少量完整方案。本文把等待比较的方案称为候选方案（candidate solution）；每份候选都包含切分、核心分配和顺序。第三，检查覆盖、重复安排和循环依赖。第四，在有证明依据时，用完成时间下界排除无法改善的候选。最后，完整评价剩余方案并选择输出。

图 \ref{fig:framework} 给出共同步骤；三条构造分支各自针对不同的输入关系，随后接受同样的完整方案检查。

\begin{figure}[!htbp]
\centering
\includegraphics[width=\linewidth,height=.84\textheight,keepaspectratio]{figures/fig-framework-v2.png}
\caption{三问共用的方案构造与比较步骤。各分支只在结构条件满足时采用，所形成的完整方案需检查依赖与容量，并通过执行评价比较完成时间。下界只在其证明条件成立时用于提前排除候选。}
\label{fig:framework}
\end{figure}

三问采用的局部计算各不相同。问题一计算重新分配独立分支后各核心的工作量；问题二计算生成或读取同一张量的操作分布在多少个核心；问题三计算先完成的子树留下多少字节。它们分别用于提出方案，不自动等于整个程序的Makespan。

这一组织方式使数学推导直接对应算法中的决定：无环性证明回答``数据依赖与所安排的任务顺序是否存在循环冲突''，字节公式回答``改动减少了多少边界搬运''，树排序证明回答``在指定模型内应先算哪个分支''。完成时间则由包含全部约束的评价程序确定。不同图不必采用同一个局部改动，但选择规则和评价次数上限在运行前固定。

\subsection{子图划分与执行顺序的无环性条件}

通过以下反例，我们说明仅检查每个子图内部还不足以保证任务顺序可行。原图有四个操作\(a_1,a_2,b_1,b_2\)，只存在依赖\(a_1\to b_1\)和\(b_2\to a_2\)，因此原图无环。若把两个\(a\)操作合成子图A、两个\(b\)操作合成子图B，则子图关系同时有A\(\to\)B和B\(\to\)A。两组内部即使没有环，组间也已经形成环。因此，检查任务顺序时还需要保留不同分组之间的数据依赖。

为了避免这种情况，可以为各子图建立与全部先后关系一致的编号。下面给出判断条件。

\textbf{命题3.1。} 若能给每个子图\(s\)赋一个整数\(r(s)\)，并使\(D^+(P)\)的每条边\(s\to t\)都满足\(r(s)<r(t)\)，则该图不存在有向环。

\begin{proof} 假设存在\(s_1\to\cdots\to s_m\to s_1\)，沿环逐条应用不等式将得到\(r(s_1)<r(s_1)\)，矛盾。\end{proof}

这里\(r\)只是用于证明先后关系的编号函数，不是性能评分。局部重分组时，若能够为新子图安排满足上述条件的编号，就证明了不会引入循环等待；存储容量是否足够、时间是否更短，仍需另外检查。

\subsection{完成时间下界与候选方案的排除条件}

这里要限制的量是全部操作完成所需的模拟时间，即总体任务执行时间\(M_q(P)\)，单位为时钟周期（cycle）。固定输入图\(G\)、核心数量\(K\)、问题编号\(q\)和官方配置后，如果每一份合法方案\(P\)都满足\(M_q(P)\ge L\)，就称\(L\)为该问题的\textbf{完成时间下界（lower bound on makespan）}。\(L\)的单位同样是cycle；它表示任何合法方案都至少需要这么多时钟周期，不是求解程序运行秒数的下界，也不是数据搬运字节数的下界。

先利用所有合法方案都必须满足的计算条件求一个下界：每个操作保留原有计算时长，每个核心的同一Pipe同时至多执行一个操作，且计算过程必须遵守已确认的数据依赖。暂时不计数据搬运、缓存换入换出和任务切换造成的额外等待，只会放宽这些条件，不会要求更多完成时间。

记\(W_M\)为全部Cube操作的时长之和，\(W_V\)为全部Vector操作的时长之和，单位均为cycle。对保留的计算操作之间的数据依赖图，将一条路径上各操作的计算时长相加；这些路径时长的最大值记为\(L_{\mathrm{cp}}\)，称为关键路径长度（critical path length），单位也是cycle。这里的``长度''指所需时间，不是节点数或边数。因为同一路径上的操作必须依次完成，且每类Pipe仅有\(K\)份同类执行资源，所以

\[L_{\mathrm{basic}}=\max\left\{L_{\mathrm{cp}},\left\lceil\frac{W_M}{K}\right\rceil,\left\lceil\frac{W_V}{K}\right\rceil,\max_{v\in V}d_v\right\}\le M_q(P).\]

式中每一项都以cycle计，\(d_v\)是操作\(v\)的计算时长，\(\lceil x\rceil\)表示向上取整；不等式对上述固定条件下的每份合法方案\(P\)成立。Cube与Vector是不同资源，可以同时工作，因此应分别限制其工作量，不能将\(W_M+W_V\)除以\(K\)后作为同一资源的必要时间。依赖图也只能保留有依据的计算关系；经过原始COPY的路径不能任意删除中间节点后就当成直接计算关系。

另一类下界只约束一份已经固定的候选方案\(P\)：根据这份方案的核心分配、各Pipe顺序及必需的数据传递等待，求得\(L(P)\le M_q(P)\)。\(L(P)\)仍以cycle计，但不保证对其他方案成立。若已有合法方案\(P_0\)，其完成时间为\(U=M_q(P_0)\)，且\(L(P)\ge U\)，就能断定候选\(P\)不会比\(P_0\)严格更快，可以省去对\(P\)的完整评价。这个判断不能推出所有合法方案的最优完成时间都大于\(L(P)\)。

若\(L\)是对全部合法方案成立的完成时间下界，\(U\)是同一问题中一份合法方案的完成时间，则\(L\le\mathrm{OPT}\le U\)。其中\(\mathrm{OPT}\)是全部合法方案中最小的Makespan；\(L\)、\(\mathrm{OPT}\)和\(U\)均以cycle计。在\(L>0\)时，两侧同单位的时间相除，得到无量纲的比值，并有

\[\frac{U}{\mathrm{OPT}}-1\le\frac UL-1.\]

例如，\(U/L\le1.05\)说明该方案的完成时间至多比最优完成时间长5\%。这个百分比衡量完成时间的相对偏差。反过来，比值较大可能是算法还有改进空间，也可能只是下界忽略了较多约束，不能把这部分差值都当成可以实现的收益。

\subsection{求解时间的组成与计算复杂度}

求解程序的实际运行时间与上一节的模拟时间分开计算。设程序用于解析输入、构造方案、检查方案、读写文件的时间分别为\(T_{\mathrm{parse}}\)、\(T_{\mathrm{construct}}\)、\(T_{\mathrm{check}}\)和\(T_{\mathrm{IO}}\)，第\(j\)次完整评价的实际运行时间为\(T_j\)，这些量的单位都是秒。令\(b\)为实际评价次数，按互不重复的程序阶段计时，则

\[T_{\mathrm{solve}}=T_{\mathrm{parse}}+T_{\mathrm{construct}}+T_{\mathrm{check}}+\sum_{j=1}^{b}T_j+T_{\mathrm{IO}}.\]

这个式子单独列出完整评价所用的时间，因为它会随图规模和存储事件数变化。算法复杂度另外描述基本操作次数随输入规模的增长：若操作之间保留的依赖集合为\(E_V\)，一次覆盖和拓扑检查需要\(O(|V|+|E_V|+|\mathcal S|)\)次基本操作。基本操作次数不能与秒数直接相加，也不等于程序的实测时间。三种算法的构造步骤、评价次数上限和实测求解时间在对应章节说明。

进一步减少中间表示的信息量，需要证明压缩前后能够判断相同的操作、得到相同的后续结果。仅有当前完成时间相等还不够。附录给出这一要求的形式化说明；本文没有把这种一般压缩方法列为已经实现的算法。


\section{问题一：子图对应的任务分别执行时的调度}

\subsection{任务划分的并行收益与通信开销}

问题一中，每个子图中所表示的操作构成一个Task。拆分能够把其中一部分计算交给另一个核心，但新Task之间的数据必须经过DDR；即使它们最终仍分到同一个核心，也不能保留跨Task数据。因此，算法需要判断新获得的并行计算是否值得新增的搬运和等待。

设Task \(s\)在核心 \(k\)的开始、完成时刻分别为\(A_s,F_s\)。记\(\mathrm{prev}(s)\)为同一核心上紧邻它的前一个Task，\(\operatorname{pred}(s)\)为它所依赖的Task集合。题面至少要求

\[A_s\ge F_{\mathrm{prev}(s)}+100,\]

\[A_s\ge F_t+1000,\quad t\in\operatorname{pred}(s),\ a(t)\ne a(s).\]

不存在的前驱不参与约束，时刻单位为cycle。这些不等式说明拆分会增加哪些开始条件，却不能单独算出\(F_s-A_s\)：Task执行期间还可能与其他核心争用DDR带宽。仅把操作周期相加，也不足以得到其最终持续时间。

\subsection{根据数据依赖结构构造调度方案}

算法首先识别计算图中彼此没有计算操作之间的数据依赖的部分。把有依赖的操作间连线暂时看作无向边，能够沿连线互相到达的操作构成一个连通分量（connected component）。不同连通分量可以分别安排，但若它们读取同一外部输入，搬运仍可能相互影响。

固定实现最多产生六种基础方案，分别考虑：整体分配这些独立部分；处理工作量集中的较大部分及末端汇合计算；拆分工作量过多的部分；安排读取共享输入的计算；分开能够独立计算的分叉支路；处理矩阵计算之后经过一段向量计算、再进入矩阵计算的受限链状结构。最后一类只在中间结果不被其他分支读取、并满足容量检查时采用。

在基础方案上，还可按固定条件增加三类比较：调整上述受限链的分组长度；尝试两种保持局部依赖结构的计算顺序；把符合条件的一个独立分支改分到另一核心。六份基础方案、至多一份分组长度调整方案、至多两份顺序调整方案和至多一份分支重分配方案，使在线E1评价总数不超过十次。不满足构造条件或与已有方案完全相同的候选不进入评价。

E1根据完整方案评分，用于当前输入的方案选择；主结果由对应的E0原件确认。

\begin{figure}[!htbp]
\centering
\includegraphics[width=\linewidth,height=.84\textheight,keepaspectratio,trim=0 235 0 0,clip]{figures/fig-p1-cuts-zh.png}
\caption{case\_026、5个核心的真实归约片段。左为原始划分，右为将独立支路改分到另一个核心；八个 ADD 节点的编号和依赖不变。点线端口连接裁剪范围外的输入或后继，橙色虚线为位于不同子图中的操作之间的依赖。}
\label{fig:p1-cuts}
\end{figure}

\subsection{独立计算分支的重新分配及无环性证明}

以下说明上一节最后一类改动。对原方案的子图依赖图，定义\(h(s)\)为从起点到Task \(s\)的最长依赖路径边数；没有前驱时\(h(s)=0\)。算法只处理同一核心上的Task按\(h\)严格递增排列的方案，并要求同时参与改动、具有相同\(h\)的不同Task之间没有计算操作之间的数据依赖。

在工作量较大的一个Task内，找到两条互不依赖、随后汇合的计算支路。把相关操作划为三个互不重叠的集合：\(Y\)为准备交给另一核心的独立支路，\(X\)为留在原核心的其余前部计算，\(J\)为等待两侧结果的汇合计算及其后续。要求\(X\)与\(Y\)之间没有依赖，\(J\)不能返回依赖到\(X\)或\(Y\)，与外部Task相连的边仍遵守原来的\(h\)先后关系。

原来承担该Task的核心先后执行集合\(X,J\)中的操作；接收\(Y\)的核心先执行集合\(Y\)中的操作，再执行其原有的、具有相同\(h\)值的Task。因此，程序必须确认接收核心原先存在这个Task，而且在同一次修改中不会再从接收核心向其他核心转出支路。算法根据各核心的计算工作量，估计重新分配后最大的工作量是否严格减少。只有满足此条件，才继续生成需要完整评价的方案。

上述调整同时改变数据在哪个核心生成、哪些任务必须先完成。只看两条支路互不依赖，还不能排除它们与原有任务次序产生冲突。因而需要检查整个调整后的顺序，以下命题给出足够的条件。

\textbf{命题4.1。} 在上述条件下，保留不同\(h\)之间的原顺序后，这种分支重新分配不会在加入同一核心顺序后的子图依赖图中引入有向环。

\begin{proof} 先按\(h\)的大小排列Task。对\(h\)相同的Task，将独立的\(X,Y\)放在前面，将接收核心原有的Task放在收到的\(Y\)之后，将\(J\)放在它所需的\(X,Y\)之后；若接收多个\(Y\)，再依实际次序编号。原方案中这些\(h\)相同的Task互无计算操作之间的数据依赖，\(X,Y\)也互不依赖，\(J\)没有返回边，且接收核心不同时向外转出支路。因此，在每个固定\(h\)内，数据依赖及同一个核心中的执行顺序都与这个编号一致。连接不同\(h\)的边则严格增加\(h\)。由命题3.1，所得图无环。\end{proof}

这个证明只说明执行顺序不会形成循环。新增子图仍可能增加复制，改变存储申请和操作等待。因此，形成完整方案后仍需检查全部操作覆盖与顺序，再用E1比较完成时间及代码规定的COPY统计量。

\subsection{算法步骤与停止规则}

\textbf{算法1：问题一的切分与调度。}

\begin{quote}
输入：原图G、核心数量K、官方固定配置。

1. 按已固定的六类结构条件产生基础方案，删除内容完全相同的重复方案。

2. 比较需要评价的基础方案，保留当前最优的完整方案。

3. 若受限链状结构满足要求，增加一次分组长度调整并比较。

4. 按固定的两种规则分别调整局部计算顺序，检查并比较所得方案。

5. 若满足本章定义的高度与依赖条件，构造一次X、Y、J分组及分支重新分配。

6. 检查全部操作覆盖、每个子图只安排一次，以及完整顺序无环。

7. 仅按规定的完整评分接受改动，输出最佳合法方案；总E1调用不超过十次。
\end{quote}

读取、覆盖检查和拓扑检查的时间与图中节点、边数量成正比；除此之外，还需执行工作量排序和可重分配支路查找。分支重新分配只对Task数不超过320的方案尝试，结构检查设置2秒软预算：程序在规定检查位置发现超时后退出，并不保证在任意指令处都能立即中止。分支重新分配至多增加一次E1评价。调整分组长度的过程另有编译和子进程时间限制，均计入总求解时间。

\subsection{实验结果与执行时间分析}

固定版本在100张图、每图1～5个核心的500种组合上均得到合法结果。下表分别列出题目规定的平均加速比，以及使用算法输出方案计算的平均加速比；两行在一个核心时采用不同的方案。

\begin{table}[!htbp]\centering\small
\begin{tabularx}{\linewidth}{@{}Xrrrrr@{}}\toprule
核心数量 & 1 & 2 & 3 & 4 & 5 \\
\midrule
官方口径平均加速比 & 1.0000 & 1.9516 & 2.7514 & 3.4786 & 4.0553 \\
优化方案 $\operatorname{mean}(B/M)$ & 1.0021 & 1.9516 & 2.7514 & 3.4786 & 4.0553 \\
改善/持平/退化 & 0/100/0 & 8/92/0 & 15/85/0 & 12/88/0 & 13/87/0 \\
求解时间均值 / s & 0.436 & 3.316 & 3.426 & 3.427 & 3.522 \\
\bottomrule\end{tabularx}\end{table}

\begin{figure}[!htbp]
\centering
\includegraphics[width=\linewidth,height=.84\textheight,keepaspectratio]{figures/paper-v2/fig-p1-results.pdf}
\caption{问题一的完整版本比较。上：每个核心数量的 100 图平均加速比；下：全部500种图与核心数量组合的时间与搬运量配对变化。核心数量用点形区分，横轴采用对称对数。}
\label{fig:p1-results}
\end{figure}

5个核心时，平均加速比由前版4.0259提高到4.0553，相对增幅为0.729\%。分支重新分配在全部500种组合中被选用36次，其余保留此前方案。上述版本比较同时包含其他规则的变化，尚未单独测得分支重新分配对全部改善的贡献。

5个核心时的额外DDR搬运总量为1053684536 B，比前版增加约0.389\%。因此这次版本改进缩短了完成时间，但没有同时减少总搬运。求解时间的95\%分位数为16.623秒，最大为25.552秒。该批次在共享主机上用四个工作进程执行，没有清空操作系统缓存；这些时间不代表独占机器上的确定延迟。

全部测试包括500次新求解进程和1130次在线E1评价。48种组合进行了新的E0评价；另外452种组合在方案文件、原图、核心数量、配置和评价规则逐项一致后，引用已有的E0结果。旧E0的运行时间没有记为新测值。

图 \ref{fig:p1-overlap} 将修改前后的完整方案放在相同时间尺度上比较，可直接观察各核心上最后一个操作的完成时刻，以及不同Pipe上各操作的执行区间。

\begin{figure}[!htbp]\centering
\includegraphics[width=\linewidth]{figures/paper-v2/fig-p1-overlap-old.pdf}
\par\vspace{4pt}
\includegraphics[width=\linewidth]{figures/paper-v2/fig-p1-overlap-new.pdf}
\caption{case\_026、5个核心 两种切分的执行时间线。上：修改前方案，$M=44{,}114$ cycles，额外搬运211,680 B；下：分支重新分配后，$M=42{,}014$ cycles，额外搬运375,532 B。每组分别显示Task持续区间及四条Pipe的操作区间，横轴范围相同，数据来自官方模拟事件。}
\label{fig:p1-overlap}
\end{figure}

从这组配对结果可见，把独立计算交给另一个核心能够减少串行工作，但也要付出新的搬运和等待。问题二允许同一个核心保留不同子图间传递的数据，其边界搬运量需要按新的Task组成方式计算。


\section{问题二：联合安排计算位置与数据搬运}

\subsection{同一核心内的数据复用}

问题二把分配到同一个核心的各子图所对应的操作合并为一个执行任务（Task）。先处理的子图中操作所生成的数据，可以留在这个核心的L1/UB中，供后面子图中的操作使用。如果生成和读取数据的操作分配到了不同核心，数据仍须通过COPY\_OUT写入DDR，再由目标核心的COPY\_IN读入；目标读取须等待源COPY\_OUT完成后500 cycles。此时，把每个计算操作单独划成子图，也不意味着为每个操作建立独立Task。提交文件中的子图顺序会影响官方程序在这个核心中安排操作的先后次序。

这种变化使分配与排序更加相关。把读取同一张量的操作安排到同一个核心，可能只需搬入一次；但如果两次使用相距较远，这块数据就要保存较长时间，期间可能妨碍其他数据分配空间。把这些操作分到不同核心虽然增加搬运，却可能让计算更早完成。本文分别估计可安排计算的时间和边界复制字节数，再比较完整方案。

\subsection{利用执行资源空闲区间安排计算}

算法记录每个核心上Cube和Vector已经安排的计算时间区间。安排下一个操作时，先根据保留的计算操作之间的数据依赖估计它最早能够开始的时刻，再寻找同类执行单元上足够长的空闲区间。这样，操作可以插到两次已安排计算之间，不必总是排到最后。

设\(v\)的估计最早开始时刻为\(r_v\)，计算时长为\(d_v\)；\(\mathcal I_{k,p_v}\)是核心 \(k\)的Pipe \(p_v\)已经占用的时间区间的并集。选取最早的\(s_v\)使

\[s_v\ge r_v,\qquad [s_v,s_v+d_v)\cap\mathcal I_{k,p_v}=\varnothing.\]

区间采用左闭右开表示，因此一次操作恰好在另一次完成时开始，不算同时占用资源。两条计算支路随后汇合时，算法联合比较它们的核心分配，至多检查\(K^2\)个核心组合，避免只看一侧较早完成而忽略另一侧。

这个估计没有完整模拟DDR带宽分配、缓存换入换出和存储空间分配次序。因此，\(s_v\)用于候选构造，正式执行时刻由完整模拟确定；该估计不具备完成时间上界或下界的保证。算法最终用单操作子图及其列表顺序提交结果，由官方程序决定实际开始时刻。

同一个核心保留数据与在不同核心之间复制数据的区别见图 \ref{fig:progression}。后续按这个区别计算边界复制量。

\subsection{基于张量关联操作集合的复制量计算}

先看一个简单例子。计算操作\(u\)生成8 B的张量\(t\)，操作\(v_1,v_2\)都读取\(t\)。如果两个读取操作分到同一个远端核心，就应按该核心需要的数据副本统计，而不能把两条读取关系视为互不相关的两次向另一个核心传递数据的需求。若只把\(v_1\)移回\(u\)所在的核心，\(v_2\)仍需远端副本，相关边界复制量未必下降；将两个读取操作一起移回后，才消除这份远端需求。

为表示这种多操作共同使用一块数据的关系，引入超边（hyperedge）：它是一个顶点集合，可以同时包含两个以上顶点；由顶点与这些集合组成的结构称为超图（hypergraph）{[}3{]}{[}5{]}。本文以计算操作为顶点，对中间张量\(t\)建立集合\(e_t=\{u,v_1,v_2,\ldots\}\)，其中包含生成\(t\)的计算操作和所有读取\(t\)的计算操作。超边只表示共同关联，不表示执行方向；原来的有向依赖仍单独保留。

应用以下公式前要求：全部非COPY计算操作已完整纳入模型，原图不同张量标识不通过逻辑标识合并，每个中间张量至多由一个原始计算操作生成，张量字节数为非负整数。这里的``生成''指计算得到数据，不把后来COPY得到的副本当成另一次原始计算生成。

给定核心分配\(a\)，定义\(\lambda_e(a)\)为超边\(e\)中各操作所在的不同核心数量。令\(w_e\)表示新增一个相关核心所增加的边界复制字节数，\(C\)表示与分配无关的固定复制量，\(\mathcal H\)为全部需计入的超边集合。在官方Step2插入容量不足所需的缓存换出、换入之前，复制量为

\[D_{\mathrm{pre2}}(a)=C+\sum_{e\in\mathcal H}w_e\bigl(\lambda_e(a)-1\bigr).\]

式中\(D_{\mathrm{pre2}},C,w_e\)的单位都是B，\(\lambda_e\)是无量纲整数。下标pre2表示这项复制量在官方Step2插入缓存换出、换入之前统计。

对DDR中的外部输入，只把读取该数据的计算操作放入相应集合。每增加一个需要该输入的核心，就增加\(b_t\) B的搬入量，所以\(w_e=b_t\)，第一次输入复制计入\(C\)。对由计算产生的中间张量，集合还包含生成它的操作；官方边界构造对每个远端核心计入输出、输入一对复制，所以\(w_e=2b_t\)。最终输出等与分配无关的部分也计入\(C\)。两类数据对应的权重分别为\(b_t\)与\(2b_t\)。

将上述公式用于8 B示例，当\(u\)在核心 0、两个读取操作在核心 1时，\(\lambda_e=2\)，该中间张量贡献16 B的边界复制。只移回一个读取操作，仍贡献16 B；两个都移回后，\(\lambda_e=1\)，该项变为0。

\begin{figure}[!htbp]
\centering
\includegraphics[width=\linewidth,height=.84\textheight,keepaspectratio]{figures/fig-hypercut-v2.png}
\caption{同一张量的计算操作分配。图中$u$生成$t$，$v_1,v_2$读取$t$；$\lambda_t$表示这三个操作所在的不同核心数量。16 B、16 B、0 B仅统计该张量的边界复制，线条只表示集合关联。}
\label{fig:hypercut}
\end{figure}

还可以只更新受移动影响的超边。设一组操作\(A\)全部从核心 \(k\)移至不同的核心 \(l\)。对与\(A\)相交的超边\(e\)，令\(n_{e,k},n_{e,l}\)为移动前两个核心中属于\(e\)的操作数量。若目标核心原先没有相关操作，则增加一次权重；若源核心的相关操作全部被移走，则减少一次权重。因此

\[\Delta_e=w_e\left(\mathbf1\{n_{e,l}=0\}-\mathbf1\{n_{e,k}=|e\cap A|\}\right).\]

\(\mathbf1\{\cdot\}\)在括号条件成立时为1，否则为0；没有与\(A\)相交的超边变化为0。这个式子能够计算字节数变化，却仍未包括缓存容量不足引发的额外搬运，也没有描述计算与搬运能否同时进行。

\subsection{计算工作量约束下的局部最小割方法}

算法在上一节的时间区间构造结果上，选择两个核心及最多16条计算链组成局部调整范围。计算链指按依赖前后相接的一串计算操作；范围之外的操作以及必须保持位置的链不允许更换核心。其余操作在所选两个核心之间选择位置。

对于固定范围与固定操作，目标是减少这一范围影响的边界复制量。借助辅助节点，可以把两种核心选择表示为源点与汇点之间的最小割问题（minimum source--target cut）：在一个带非负边权的有向网络中，将指定的源点和汇点分到两侧，并使从源点一侧指向汇点一侧的边的权重和最小。辅助节点把``某条超边是否同时涉及两侧''转换为相应边权，因此求得的是当前局部复制字节目标的精确最小值。

只限制字节数可能把过多计算集中到同一核心。实现因而还规定每个核心、每类计算Pipe的累计计算时长上限。若新的分配超过上限，就固定部分移动链回到原核心，再求解缩小后的局部问题。原分配始终保留为可行选择；若局部有\(r\)条可调整链，\(r\le16\)，至多求解\(r+1\)次。接受前重新计算工作量与复制量，两项检查都通过才保留改动。

\textbf{命题5.1。} 若每次固定操作都与原分配一致，每次最小割得到当前局部问题的精确最优值，且接受前核对计算时长上限，则被接受方案的局部\(D_{\mathrm{pre2}}\)不大于原分配。

\begin{proof} 原分配在每次增加固定操作后仍可行，因此每次局部最优值都不大于原分配值。最终只接受满足计算时长上限且重新核算复制量不增加的方案，结论成立。\end{proof}

命题保证的是固定局部范围内复制量不增加。固定部分操作会缩小可选分配集合，而Makespan还取决于操作次序和执行资源。算法保留调整前的完整方案，让最终评价决定是否采用新的分配。

\subsection{完整算法及计算量}

\textbf{算法2：问题二的切分与调度。}

\begin{quote}
输入：原图G、核心数量K、官方固定配置。

1. 按图结构建立基础方案P0。

2. 在各核心的计算空闲区间中安排操作，得到方案P1；条件不满足时保留P0。

3. 对P1中选定的两个核心和最多16条链，按边界复制量求局部最小割。

4. 检查各核心各Pipe的计算时长上限；必要时固定部分链回到原位置后重算。

5. 只有Step2之前的复制量严格减少，才根据新分配重新安排顺序，得到P2。

6. 保留未修改的P1，按P0、P1、P2的构造顺序去掉完全相同的方案。

7. 完整评价至多三份方案，先比较Makespan，再按代码规定比较新增COPY；同分保留较早方案。
\end{quote}

这里\(P_0,P_1,P_2\)是本次求解中产生的方案编号，不是问题编号。当局部复制量没有减少时，不生成\(P_2\)，保留\(P_1\)参与最终比较。

设图有\(n\)个计算操作、\(m\)条用于构造的计算操作之间的数据依赖，查找空闲区间累计扫描量为\(Q\)。若共调整\(R\)个局部范围，第\(j\)个范围的辅助网络有\(n_j\)个节点、\(m_j\)条边，含\(w_j\le16\)条可调整链，一次最小割耗时为\(T_{\mathrm{cut}}(n_j,m_j)\)，则该部分总耗时不超过

\[\sum_{j=1}^{R}(w_j+1)T_{\mathrm{cut}}(n_j,m_j).\]

此外还需计入依赖识别、排序、空闲区间扫描\(Q\)及汇合附近至多\(K^2\)个分配比较。局部最小割不调用完整评价程序，完整在线评分最多三次。全部步骤的时间均计入求解时间。

\subsection{实验结果与搬运量分析}

固定版本在全部500种图与核心数量组合上得到合法结果。官方主曲线保留使用一个核心时的加速比为1，同时单独列出算法改善仅使用一个核心的方案后的比值。

\begin{table}[!htbp]\centering\small
\begin{tabularx}{\linewidth}{@{}Xrrrrr@{}}\toprule
核心数量 & 1 & 2 & 3 & 4 & 5 \\
\midrule
官方口径平均加速比 & 1.0000 & 2.3167 & 3.2356 & 3.9711 & 4.5498 \\
优化方案 $\operatorname{mean}(B/M)$ & 1.2061 & 2.3167 & 3.2356 & 3.9711 & 4.5498 \\
改善/持平/退化 & 59/41/0 & 64/36/0 & 54/46/0 & 49/51/0 & 42/58/0 \\
\bottomrule\end{tabularx}\end{table}

与前一完整版本相比，268种组合的Makespan缩短，232种持平，没有变慢；69张图至少在一个核心数量下得到改善。该比较评价的是完整算法在给定数据集上的表现；各局部方法的独立作用尚未由这组比较区分。

全部组合的额外DDR总量从3043860416 B增至4159738560 B，增加约36.66\%。其中55种组合的时间与搬运同时减少，199种时间减少而搬运增加，14种时间减少、搬运不变，232种两项均不变。算法选择以完成时间为主，因此会接受这类以更多搬运换取更早完成的方案。

平均求解时间为4.029秒，95\%分位数为18.841秒，最大40.187秒。独立E0评价的平均时间为1.410秒，95\%分位数为6.592秒，最大13.980秒。这批结果在共享Mac上用一个工作进程取得，整批约46.63分钟；每次求解新建进程，操作系统缓存保持原有状态。

\begin{figure}[!htbp]
\centering
\includegraphics[width=\linewidth,height=.84\textheight,keepaspectratio]{figures/paper-v2/fig-p2-results.pdf}
\caption{问题二的完整版本比较。上：按题目要求计算的平均加速比；下：500种组合配对结果。部分时间改善伴随额外搬运增加，两项指标分别报告。}
\label{fig:p2-results}
\end{figure}

\subsection{减少搬运量后完成时间增加的反例}

为检验只按容量和搬运量构造方案是否足够，已有实验对另一种固定构造C04评价了六个预先选择的5个核心用例。C04是该实验的代码标识。六份方案都合法，且官方Step2未插入因容量不足产生的缓存换出、换入，但完成时间全部长于本章主算法。

{\def\LTcaptype{table} % do not increment counter
\begin{table}[!htbp]\centering\small
\begin{tabular}{@{}lrrr@{}}
\toprule\noalign{}
case & 主算法 / cycles & C04 / cycles & 相对变慢 \\
\midrule\noalign{}


005 & 33515 & 38209 & 14.01\% \\
016 & 2091375 & 2399125 & 14.72\% \\
019 & 16247 & 28514 & 75.50\% \\
025 & 942938 & 975658 & 3.47\% \\
039 & 449653 & 494942 & 10.07\% \\
075 & 343930 & 452120 & 31.46\% \\
\bottomrule\end{tabular}
\end{table}
}

其中case\_019、case\_039、case\_075的COPY字节减少，完成时间却增加。由这三个反例可知，减少搬运量且消除容量不足引起的换入换出，仍不足以保证完成时间缩短。具体原因还需结合执行事件分析。即使计算操作的顺序不变，搬运操作的顺序、不同核心之间传递数据所需的等待及DDR请求同时发生的情况仍可能改变。

我们在图 \ref{fig:p2-counterexample} 中对照两份方案的官方执行事件：搬运量较小的方案，最后一个操作反而完成得更晚。

\begin{figure}[!htbp]\centering
\includegraphics[width=\linewidth]{figures/paper-v2/fig-p2-counterexample-control.pdf}
\par\vspace{4pt}
\includegraphics[width=\linewidth]{figures/paper-v2/fig-p2-counterexample-c04.pdf}
\caption{case\_019、5个核心的官方模拟的执行时间线。上：主算法，$M=16{,}247$ cycles；下：C04，$M=28{,}514$ cycles。每组分别显示Task和四条Pipe，共用横轴范围。C04的搬运量较少，且没有容量不足引起的缓存换入换出，但完成时间更长。}
\label{fig:p2-counterexample}
\end{figure}

本章用边界复制公式选择调整方向，再用完整模拟比较执行时间。问题三增加共享只读Cache后，还要判断读取请求到达时数据是否在Cache中，单看操作分到哪个核心更加不足。


\section{问题三：安排计算顺序并利用共享只读Cache}

\subsection{共享只读缓存的访问规则}

问题三保留问题二的Task组成和分配到不同核心的操作之间的数据依赖，增加全部核心共享的只读Cache。按官方代码，操作类型为COPY\_IN、能够确定搬入数据的逻辑标识且搬运字节数大于0时，操作开始就查询Cache。程序使用这一逻辑标识查找数据。若数据存在，则为命中（cache hit），从Cache读取；否则为未命中（cache miss），从DDR读取，并在传输完成后把数据放入Cache。FIFO规则按插入时间选择需要移除的对象，命中不刷新该对象的次序。单份数据大于Cache容量时，不将它插入Cache。

因此，前一次DDR传输尚未完成时，后续请求不能仅因数据正在传输就算作命中。即使命中，也仍须等待原有数据依赖及不同核心之间传递数据时的500 cycles等待条件满足，并占用Cache读取带宽。改变计算顺序会改变读取请求出现的时刻，而总命中字节数并不能完整描述这种时间关系。

本章算法在基础分配上，按输入结构尝试共享输入的访问顺序、符合条件的注意力计算顺序，以及下文定义的树状计算顺序。本章完整算法在代码中记为Forest，其中树排序仅用于满足本节条件的输入。

例如，case\_021、3个核心的编译后张量 t1000000001 大小为 4096 B。官方事件给出以下生命周期；事件时刻为 cycles，表中间隔不表示等长处理阶段。

\begin{table}[!htbp]\centering
\caption{同一张量的 Cache 事件}\label{tab:cache-events}
\begin{tabularx}{\linewidth}{rX}\toprule
时刻 / cycle & 事件 \\\midrule
0 & COPY\_IN开始时未命中 \\
207 & DDR COPY\_IN 完成，张量放入Cache，记下插入顺序 \\
4,956 & 另一个核心读命中，插入次序不变 \\
323,850 & 插入张量 t1000000487 时，该最早插入的项被移除 \\\bottomrule
\end{tabularx}\end{table}

\subsection{子树处理顺序对结果存储峰值的影响}

考虑两个互不依赖的计算分支A、B。父操作必须等两者完成，才能读取它们的结果。A先完成时，它留下的结果不能立刻释放，因为父操作还要等B；这部分结果要与B计算期间的数据同时保留。反过来先计算B，保留的数据量可能不同。

为了单独分析这个影响，先固定一棵树状计算结构：每个内部操作只输出一个张量，该结果只被唯一的父操作读取，其他树或分支不使用它。每次完整处理一棵子树后才处理下一棵，子树之间不交错执行。对子树\(i\)定义\(h_i\)为处理它时内部计算结果占用的最大总字节数，\(r_i\)为子树完成后必须保留的输出字节数。这里统计的是此模型内的结果数据，不包括外部输入，也不分别区分L1和UB。

例如，A的\((h_A,r_A)=(8,2)\) KiB，B的\((h_B,r_B)=(7,5)\) KiB，父输出为3 KiB。先A后B时，三个需要检查的量为8、\(2+7=9\)、\(2+5+3=10\)，峰值为10 KiB。先B后A时则为7、\(5+8=13\)、\(5+2+3=10\)，峰值为13 KiB。最后一项保留了父输出先分配、子结果随后释放的规则。上述数值为用于解释公式的设定示例。

图 \ref{fig:forest-example} 按这三个处理步骤分解所占空间，保留的数据与正在处理的子树分别着色。

\begin{figure}[!htbp]
\centering
\includegraphics[width=\linewidth,height=.84\textheight,keepaspectratio]{figures/paper-v2/fig-forest-example.pdf}
\caption{两种子树处理顺序的内部结果占用。左：先A后B，峰值10 KiB；右：先B后A，峰值13 KiB。横轴为处理步骤，不是模拟时钟；数字为第6章的公式示例。}
\label{fig:forest-example}
\end{figure}

\subsection{非交错子树处理的最优顺序及证明}

设同一父操作有若干子树，\(\pi=(\pi_1,\ldots,\pi_m)\)为它们的处理顺序，父输出大小为\(r_{\mathrm{parent}}\)。当处理第\(j\)棵子树时，先前子树的输出共占用\(\sum_{\ell<j}r_{\pi_\ell}\)，当前子树最多再占用\(h_{\pi_j}\)。父操作开始分配输出时，各子树输出仍需保留。因此峰值为

\[H(\pi)=\max\left\{\max_j\left(\sum_{\ell<j}r_{\pi_\ell}+h_{\pi_j}\right),\ \sum_i r_i+r_{\mathrm{parent}}\right\}.\]

所有项的单位均为B。该式允许比较排序，却不直接包含执行时间。

\textbf{命题6.1。} 若\(h_i\ge r_i\ge0\)，在上述不交错处理子树的模型中，按\(h_i-r_i\)从大到小排序能够使\(H\)最小。

\begin{proof} 取相邻两棵子树A、B，假设\(h_A-r_A\ge h_B-r_B\)。暂时不计两者之前已经保留的共同结果，A后B与B后A的最大占用分别为

\[H_{AB}=\max(h_A,r_A+h_B),\qquad H_{BA}=\max(h_B,r_B+h_A).\]

由假设可得\(r_A+h_B\le r_B+h_A\)，又因\(r_B\ge0\)而有\(h_A\le r_B+h_A\)，故\(H_{AB}\le H_{BA}\)。之前保留的数据只会给两种次序各加同一个常数；两棵子树都完成后，各自留下的结果之和同为\(r_A+r_B\)，后续计算及父输出分配不受交换影响。不断把次序不符合要求的相邻子树交换，最大占用不会增加，最终得到按\(h_i-r_i\)递减的顺序。对子树自下而上应用这一结论即可。\end{proof}

这种通过交换相邻分支证明存储次序的方法，与表达式树代码生成中的经典研究相关{[}4{]}。针对本题按字节给定大小的结果数据，我们以上述相邻交换论证证明相应结论。

算法先确认每个内部输出只被父操作读取、原COPY只位于允许的边界，并且依赖中的数据大小完整，再自下而上计算\(h_i,r_i\)。按所求顺序先处理子节点、后处理父节点，保持原来的核心分配，用单操作子图列表提交新顺序。

该最优性结论适用于前述单棵树、内部结果、子树不交错处理的模型。完整模拟还计入外部输入、不同树的最终结果、L1与UB的独立容量以及新建数据副本；算法据此进一步评价新顺序的缓存换入换出和完成时间。

\subsection{固定候选方案的完成时间下界}

当一份候选的单操作子图、核心分配和各核心各Pipe顺序已经固定时，可以列出每次成功执行都必须满足的时间不等式。对必要的先后关系\(u\to v\)，写为

\[s_v\ge s_u+d_u+\delta_{uv}.\]

\(s_u,s_v\)是开始时刻，\(d_u\)是计算时长，\(\delta_{uv}\)只计入能够证明不可避免的等待，单位均为cycle。同一对操作之间若有多个条件，应取最大的必要等待，不把重复条件相加。

将这些关系组成有向图，其最长路径长度\(L(P)\)不大于固定候选的实际完成时间。如果当前已成功评价方案的完成时间为\(U\)，且\(L(P)\ge U\)，该候选就不能严格改善，可以省去一次E0评价。

上述下界由每次合法执行均满足的时间不等式推得。放宽约束后，任意一个贪心调度结果仅是一份可行安排的完成时间；只有求得该放宽问题的最优值，或另证必要不等式，才能据此取得下界。

\subsection{算法步骤与复杂度}

\textbf{算法3：问题三的切分与调度。}

\begin{quote}
输入：原图G、核心数量K、官方固定配置。

1. 按结构与计算空闲区间构造基础方案，并进行评价。

2. 对满足条件的输入，尝试连续安排共享输入访问或调整注意力计算的次序。

3. 在剩余评价次数允许时，识别本章定义的树状计算，按子树峰值减去输出字节数的差值排序，产生方案。

4. 跳过与已有方案完全相同的结果；对满足证明条件的候选计算完成时间下界。

5. 若下界不小于当前最佳完成时间，则跳过完整评价。

6. 否则调用E0，只有Makespan严格缩短才替换当前方案。

7. 一次求解共用至多三次E0预算，输出最佳成功方案及调用记录。
\end{quote}

全部候选共用三次E0调用预算。若前面的候选已使用全部次数，后面的树排序就不会追加评价。结构不符合、方案重复、下界无法改善、合法性不通过和评价后未改善分别作为候选停止原因记录。

给定含\(n\)个节点、\(m\)条依赖的树状结构，识别及按子先父后遍历需要\(O(n+m)\)时间。每个父节点对孩子排序，合计需要\(O(\sum_v\deg(v)\log\deg(v))\)时间，其中\(\deg(v)\)为其孩子数。候选约束图含\(m'\)条必要边时，最长路径可在\(O(n+m')\)时间内求得。完整算法的计算量还包括基础分配和至多三次官方评价。

\subsection{启用和关闭共享缓存的配对实验}

Forest固定版本完成全部500种图与核心数量组合，均得到官方成功结果，500次求解共调用973次在线E0，没有重试。与前一完整策略相比，24种组合的Makespan缩短、476种持平，没有变慢。额外COPY和因容量不足产生的搬运分别有18种减少、479种持平、3种增加，两种统计各净减少413450240 B。其中，因容量不足产生的搬运属于额外COPY的一部分。

为比较Cache作用，对每份最终方案分别运行问题二、三的官方评价。下表``无L2''是同一份P3方案关闭共享Cache后的结果，并不是问题二主算法生成的另一份方案。

{\def\LTcaptype{table} % do not increment counter
\begin{table}[!htbp]\centering\small
\begin{tabular}{@{}
  >{\raggedright\arraybackslash}p{(\linewidth - 8\tabcolsep) * \real{0.1579}}
  >{\raggedleft\arraybackslash}p{(\linewidth - 8\tabcolsep) * \real{0.2105}}
  >{\raggedleft\arraybackslash}p{(\linewidth - 8\tabcolsep) * \real{0.2105}}
  >{\raggedleft\arraybackslash}p{(\linewidth - 8\tabcolsep) * \real{0.2105}}
  >{\raggedleft\arraybackslash}p{(\linewidth - 8\tabcolsep) * \real{0.2105}}@{}}
\toprule\noalign{}
\begin{minipage}[b]{\linewidth}\raggedright
核心数量
\end{minipage} & \begin{minipage}[b]{\linewidth}\raggedleft
无 L2：\(\operatorname{mean}(B/M_2)\)
\end{minipage} & \begin{minipage}[b]{\linewidth}\raggedleft
Cache：\(\operatorname{mean}(B/M_3)\)
\end{minipage} & \begin{minipage}[b]{\linewidth}\raggedleft
\(\operatorname{mean}(M_2/M_3)\)
\end{minipage} & \begin{minipage}[b]{\linewidth}\raggedleft
逐图字节命中率均值
\end{minipage} \\
\midrule\noalign{}


1 & 1.1971 & 1.1998 & 1.0023 & 6.47\% \\
2 & 2.2949 & 2.3120 & 1.0079 & 16.25\% \\
3 & 3.1745 & 3.2587 & 1.0314 & 25.96\% \\
4 & 3.9211 & 4.0915 & 1.0512 & 24.53\% \\
5 & 4.4763 & 4.7576 & 1.0829 & 31.36\% \\
\bottomrule\end{tabular}
\end{table}
}

5个核心的逐图Cache加速比平均为1.0829。这个数值按每张图的\(M_2/M_3\)计算后取平均。按所有访问字节汇总后的命中率为47.45\%，逐图命中率算术平均则为31.36\%；前一种统计让访问字节较多的图占更大权重。

平均求解时间为4.401秒，中位数1.616秒，95\%分位数20.009秒，最大51.335秒。这些时间包含求解过程中调用E0的时间，来自共享主机四个工作进程的原始批次；

\begin{figure}[!htbp]
\centering
\includegraphics[width=\linewidth,height=.84\textheight,keepaspectratio]{figures/paper-v2/fig-p3-results.pdf}
\caption{问题三的保持方案不变的配对评价。上：无 L2 与只读 Cache 的平均基线加速比；下：逐图$M_2/M_3$的均值，每种核心数量包含100对结果。}
\label{fig:p3-results}
\end{figure}

\begin{figure}[!htbp]
\centering
\includegraphics[width=\linewidth,height=.84\textheight,keepaspectratio]{figures/paper-v2/fig-cache-gain.pdf}
\caption{500 个保持方案不变的配对结果中的字节命中率与 Cache 收益。保留低于 1 的负收益点；100 个用例的不同核心数量结果相关，不能当作 500 个独立随机样本。}
\label{fig:cache-gain}
\end{figure}

\subsection{附加比值条件对方案选择的影响}

已有后续实验R9F对5个核心候选增加了三个同时满足的接受条件：启用Cache后的\(M_3\)严格下降，关闭Cache后的\(M_2\)不增加，且比值\(G=M_2/M_3\)不下降。R9F是该实验的代码标识，这三个条件不是官方评分要求。

100张5个核心图中，96张不符合该构造条件，1张由候选下界排除，另外3张被上述联合条件拒绝，没有方案被接受。最终500份输出文件均与Forest一致。然而，被拒绝的三份候选实际上同时缩短了\(M_2\)和\(M_3\)：

{\def\LTcaptype{table} % do not increment counter
\begin{table}[!htbp]\centering\small
\begin{tabular}{@{}lrrrr@{}}
\toprule\noalign{}
case/K & 原 \(M_2\) & 候选 \(M_2\) & 原 \(M_3\) & 候选 \(M_3\) \\
\midrule\noalign{}


064/5 & 7029 & 6899 & 6981 & 6877 \\
068/5 & 131631 & 98913 & 116345 & 97971 \\
088/5 & 97305 & 66729 & 85789 & 66224 \\
\bottomrule\end{tabular}
\end{table}
}

因为完成时间为正，保持\(G\)不下降等价于

\[M_{2,\mathrm{new}}M_{3,\mathrm{old}}\ge M_{2,\mathrm{old}}M_{3,\mathrm{new}}.\]

当无Cache的时间改善得更快时，这个比值可能下降，即使启用Cache后的时间也明显缩短。在这三份候选中，比值条件拒绝了Makespan更短的方案，因而妨碍了主要目标的改善。

本节将R9F作为附加选择条件的对照实验，主表报告Forest的完整结果。R9F分两批执行400和100种组合，共使用982次在线E0，仅用于说明这项额外选择条件的实际影响。

在限定的树状计算模型中，所证排序使内部结果的存储峰值最小；在完整执行中，算法则根据官方模拟的完成时间选择方案。


\section{实验结果与分析}

\subsection{输入、对照和统计方法}

实验采用题面给出的全部100张原始计算图，核心数量分别取1～5，每题共500种输入组合。三题各使用一个固定的完整求解器版本，方案构造条件和评价次数上限在运行前确定。附录C列出各题固定程序版本；\texttt{data/source-receipt.json}记录输入、方案和评价结果的来源，\texttt{data/all-results.csv}汇总逐例数值。

全文使用三类比较。第一类用官方仅使用一个核心时的基线除以当前方案的完成时间，得到题目要求的加速比。第二类比较相同题目、相同图和相同核心数量下两个完整算法版本，观察方法修改后的总体变化。第三类保持问题三的方案不变，只切换共享只读Cache，比较能否从共享只读缓存读取数据对同一方案完成时间的影响。三类比较分别评价相对单核心基线的收益、算法版本变化的收益及共享缓存的收益。

所有结果表格和数据图均由同一份1500行结果表生成。导出时检查每题100图与五种核心数量是否完整、基线是否对应原图、问题三的两份结果是否使用同一方案，再重新计算平均值。

\begin{figure}[!htbp]
\centering
\includegraphics[width=\linewidth,height=.84\textheight,keepaspectratio]{figures/paper-v2/fig-dataset.pdf}
\caption{官方 100 张计算图的结构统计。上：非 COPY 操作数与张量大小总和；下：按矩阵计算的时长占两类计算总时长的比例排序的用例分布。}
\label{fig:dataset}
\end{figure}

\subsection{三问的主要结果}

下表报告平均加速比。先对每张图计算``官方单核心基线的完成时间／算法方案的完成时间''，再对100张图取算术平均；三问使用同一组官方基线。为了展示算法还能怎样改变仅使用一个核心的方案，表内保留实际使用一个核心时的比值；问题一、二要求的主曲线仍按题面在使用一个核心时取1。

{\def\LTcaptype{table} % do not increment counter
\begin{table}[!htbp]\centering\small
\begin{tabular}{@{}lrrrr@{}}
\toprule\noalign{}
核心数量 & P1优化方案 & P2优化方案 & P3优化方案 & P3同计划无L2 \\
\midrule\noalign{}


1 & 1.0021 & 1.2061 & 1.1998 & 1.1971 \\
2 & 1.9516 & 2.3167 & 2.3120 & 2.2949 \\
3 & 2.7514 & 3.2356 & 3.2587 & 3.1745 \\
4 & 3.4786 & 3.9711 & 4.0915 & 3.9211 \\
5 & 4.0553 & 4.5498 & 4.7576 & 4.4763 \\
\bottomrule\end{tabular}
\end{table}
}

问题一要求不同Task之间的数据经过DDR；问题二、三允许同一个核心保留在不同子图之间传递的数据，但仍受带宽、容量和不同核心之间传递数据所需的等待限制。三问之间的比较同时涉及硬件规定和算法的变化。

``P3同计划无L2''一列使用问题三算法生成的方案；``P2优化方案''由问题二算法单独求得，两者不是同一个调度结果。5个核心时，P3两列平均基线加速比之比约为1.0628，而逐图Cache加速比的均值是1.0829。差异来自先求比值还是先求平均的顺序，本文以第2章定义的逐图比较为准。

\subsection{求解时间与评价调用次数}

求解时间包含输入读取、构造与检查、在线方案评价及结果写出。问题一、二在求解过程中使用E1、E2，最终E0复评另计；问题三在求解过程中直接使用E0，其时间已经包含在求解总时间内。

{\def\LTcaptype{table} % do not increment counter
\begin{table}[!htbp]\centering\small
\begin{tabular}{@{}
  >{\raggedright\arraybackslash}p{(\linewidth - 10\tabcolsep) * \real{0.1364}}
  >{\raggedleft\arraybackslash}p{(\linewidth - 10\tabcolsep) * \real{0.1818}}
  >{\raggedleft\arraybackslash}p{(\linewidth - 10\tabcolsep) * \real{0.1818}}
  >{\raggedleft\arraybackslash}p{(\linewidth - 10\tabcolsep) * \real{0.1818}}
  >{\raggedleft\arraybackslash}p{(\linewidth - 10\tabcolsep) * \real{0.1818}}
  >{\raggedright\arraybackslash}p{(\linewidth - 10\tabcolsep) * \real{0.1364}}@{}}
\toprule\noalign{}
\begin{minipage}[b]{\linewidth}\raggedright
方法与样本范围
\end{minipage} & \begin{minipage}[b]{\linewidth}\raggedleft
在线评价总次数
\end{minipage} & \begin{minipage}[b]{\linewidth}\raggedleft
求解均值 / s
\end{minipage} & \begin{minipage}[b]{\linewidth}\raggedleft
P95 / s
\end{minipage} & \begin{minipage}[b]{\linewidth}\raggedleft
最大 / s
\end{minipage} & \begin{minipage}[b]{\linewidth}\raggedright
运行条件
\end{minipage} \\
\midrule\noalign{}


P1，100图、5 核心 & 全部500种组合1130次E1 & 3.522 & 16.623 & 25.552 & 共享主机，4个工作进程 \\
P2，全部500种组合 & 1131次E2 & 4.029 & 18.841 & 40.187 & 共享主机，1个工作进程 \\
P3 Forest，全部500种组合 & 973次E0 & 4.401 & 20.009 & 51.335 & 共享主机，4个工作进程 \\
\bottomrule\end{tabular}
\end{table}
}

表中的工作进程数说明同一批次同时运行多少个求解进程。主机还可能有其他任务，因此不同批次的时间不能直接用作独占硬件的严格速度排名。问题一的时间一栏只统计100张图的5个核心结果，调用次数则覆盖全部500种组合，两者范围不同。

95\%分位数（95th percentile，P95）表示按时间排序后接近95\%位置的值。问题一原始报告采用线性插值，问题二、三采用向上取整得到的第\(\lceil0.95n\rceil\)个次序值；\(n\)是样本数量。各方法的P95按相应原始报告的定义计算。

本批次所有成功求解进程的实测时间均低于题面建议的5～10分钟。测量范围限于上述共享主机；独占硬件、清空操作系统缓存及跨平台运行尚未测试。

\begin{figure}[!htbp]
\centering
\includegraphics[width=\linewidth,height=.84\textheight,keepaspectratio]{figures/paper-v2/fig-runtime.pdf}
\caption{三种求解程序的实际经过时间累计分布，每题500次。时间包含方案构造及在线评价，横轴采用对数刻度；批次资源不同，不据此比较独占机器的速度。}
\label{fig:runtime}
\end{figure}

\subsection{算法改进的实验依据与适用范围}

完整算法版本之间的比较，能说明修改后整体结果如何变化；某个构造被选中的记录，能说明该构造确实影响了输出；同一用例的结构图与时间线，能说明该次运行发生了什么。若多个部分同时修改，仅有前后版本比较还不能分别确定每一部分贡献多少。

问题一5个核心时平均加速比提高约0.729\%，额外DDR总量增加约0.389\%。问题二有268种组合缩短完成时间，但总额外DDR增加约36.66\%。在这些输入中，完成时间缩短与搬运量增加同时出现。计算与搬运重叠所产生的具体贡献仍需逐项分析执行事件。

问题三完整版本有24种组合的完成时间缩短，同时有3种组合的额外COPY增加。C04在所测试的六个用例上均变慢；这个结论限于这六份方案。R9F中的三份候选虽然缩短了完成时间，却因附加的比值条件被拒绝。

\subsection{完成时间的下界与最优值的偏差}

如果一个时间下界\(L\)适用于该输入的全部可行方案，而当前方案完成时间为\(U\)，则\(U/L-1\)给出相对最优差距的保守上限。问题二已有固定版本的下界核对记录：在1～5个核心下，分别有81、62、32、33、23张图满足\(U\le1.05L\)。在记录所声明的条件下，这些结果距最优不超过5\%。这是逐输入的数学保证，不是统计显著性检验。

5个核心时，平均\(B/L\)约为6.108，高于算法达到的4.550。下界忽略部分通信和容量约束时，差距也会增大，因此6.108不是已经证实可以达到的平均加速比。附录进一步分析工作量与必要等待怎样共同形成下界，但没有为这些新推导重新运行全量核对；本节仍使用已有固定结果。

\begin{figure}[!htbp]
\centering
\includegraphics[width=\linewidth,height=.84\textheight,keepaspectratio]{figures/paper-v2/fig-bound-gap.pdf}
\caption{同一输入、配置与方案范围下的时间上界和下界。$L\le\mathrm{OPT}\le U$，阴影是必要界给出的上限与当前结果之间的差距，不代表一定可实现的收益。}
\label{fig:bound-gap}
\end{figure}

\subsection{实验范围与局限}

全部100张官方图都参与过研发，结果只说明给定公开数据集上的表现，不构成对未见输入的验证。同一图的五种核心数量结果相互关联，不能直接当作500个独立随机样本估计显著性。按结构而非用例编号选择算法分支，能够避免显式按编号查表，但不能消除研发过程中接触公开图的影响。

三个主版本均在给定输入上完成规定的全部组合，尚未测试更广范围的输入。方案质量由官方模拟的时钟周期衡量，求解效率由本机程序运行秒数衡量；本次未开展真实NPU设备实验。


\section{方法评价与结论}

\subsection{主要方法与适用范围}

本文共同处理计算图切分、核心分配和子图顺序三个决策。问题一把同一核心负责的两组独立计算分配到不同核心，使两组操作能够同时进行，并检查新增Task造成的等待与搬运。问题二根据生成、读取同一张量的操作分布计算边界复制量，并在计算执行单元的空闲区间中安排操作。问题三对符合条件的树状计算选择较少占用内部结果存储空间的顺序，再检查共享只读Cache条件下的完成时间。

三问的数学分析分别回答明确的问题。沿全部子图数据依赖边和同一核心上的任务顺序边严格递增的编号，证明这些先后关系不形成有向环；张量关联集合说明为什么同时移动多个读取操作可能减少复制，而逐个移动没有立即收益；相邻子树交换证明在限定模型中怎样排列能够减小存储峰值；完成时间的必要下界则排除不可能改善的固定候选。这些结论都有具体条件，不代替完整执行评价。

算法还有三方面限制。第一，共享带宽、缓存插入移除和存储空间分配会相互影响，目前的局部表示不足以完整预测全部执行行为。第二，树状计算排序等方法依赖特定输入结构；若一个结果被多个分支读取，或输入存在不满足公式要求的标识关系，就不能直接套用原证明。第三，已有实验仍发现局部更快但被附加条件拒绝的方案，表明接受规则还可以改进。后续版本的整体收益有待重新进行全量评价。

\subsection{结论}

三个固定算法均完成100张图、每图1～5个核心的官方评价。5个核心时，问题一、二、三的平均基线加速比分别为4.0553、4.5498、4.7576；保持问题三方案不变，启用Cache相对关闭Cache的逐图平均加速比为1.0829。结果来自统一输入与冻结配置，搬运增加和候选退化也一并报告。

这些结果支持本文的求解方法：先从原图依赖和数据使用关系构造少量方案，用证明检查必要条件，再以完整模拟决定性能是否改善。三问新增的硬件能力不同，但都需要明确区分数据量、执行时间和程序求解时间。后续工作将针对已出现的反例调整分配与接受规则，并在新的固定版本上验证。


\FloatBarrier
\begin{thebibliography}{9}
\bibitem{official} 第二十三届中国研究生数学建模竞赛. A题：通用神经网络处理器下的多核调度问题及官方附件[Z]. 2026.
\bibitem{graham} GRAHAM R L. Bounds for Certain Multiprocessing Anomalies[J]. Bell System Technical Journal, 1966, 45(9): 1563--1581. DOI: 10.1002/j.1538-7305.1966.tb01709.x.
\bibitem{hypergraph} ÇATALYÜREK Ü V, AYKANAT C. Hypergraph-Partitioning-Based Decomposition for Parallel Sparse-Matrix Vector Multiplication[J]. IEEE Transactions on Parallel and Distributed Systems, 1999, 10(7): 673--693. DOI: 10.1109/71.780863.
\bibitem{sethi} SETHI R, ULLMAN J D. The Generation of Optimal Code for Arithmetic Expressions[J]. Journal of the ACM, 1970, 17(4): 715--728. DOI: 10.1145/321607.321620.
\bibitem{nist} BLACK P E. Hyperedge; Hypergraph[EB/OL]. Dictionary of Algorithms and Data Structures, NIST. [2026-09-26]. \url{https://xlinux.nist.gov/dads/HTML/hyperedge.html}.
\end{thebibliography}

\appendix
\setcounter{section}{1}
\section{补充数学推导}

\subsection{数据依赖所给出的完成时间下界}

本节假设每个需要讨论的原图中间张量至多由一个原始计算操作生成，不同原始张量标识不合并成同一逻辑数据标识。取出原始计算操作，只保留有依据的先后关系，包括原来就存在的计算关系，以及同一张量的生成操作必须先于读取操作的关系。不任意缩短经过原始COPY的路径。

在任何官方成功执行中，这些计算操作及其时长都被保留，每个核心的每条Pipe也仍同时至多处理一个操作。直接取该次执行中各计算的核心和开始、完成时刻，就得到满足保留关系的一组安排。因此，在这些必要关系下求得的最小必要时间，不会大于官方完成时间。

对于分配到不同核心的计算操作，必须核对两者之间的依赖怎样在官方执行中保留。问题一要求相关前驱Task完成后，后续Task才允许开始；问题二、三要求源COPY\_OUT完成，再经过500 cycles，目标COPY\_IN才可参与就绪判断；实际开始仍需满足其余数据依赖、执行单元和存储条件。接收端若因容量不足而需要再次读入数据，还受官方Step2次序和MTE2操作队列顺序限制，不能跳过所需的基础COPY\_IN。上述必要关系均从官方成功执行所满足的条件中提取。

\subsection{指定时间区间内的工作量下界}

在B.1的计算操作之间的数据依赖图上，为操作\(v\)定义两个不包括自身的时长：\(h_v\)为它之前最长依赖路径的时长，\(t_v\)为它之后最长依赖路径的时长。无相应前驱或后继时取0。递推式为

\[h_v=\max_{u\to v}(h_u+d_u),\qquad t_v=\max_{v\to w}(d_w+t_w).\]

本节\(h_v\)以cycle计，与第6章按字节定义的子树存储量\(h_i\)是不同的局部记号。取\(a,b\ge0\)，选定一种Pipe \(p\)，考虑非空操作集合

\[S=\{v:p_v=p,\ h_v\ge a,\ t_v\ge b\}.\]

任何成功执行中，\(S\)的操作都只能完整安排在\([a,M-b]\)内：它们之前至少需要\(a\)个周期，之后至少需要\(b\)个周期。同类Pipe共有\(K\)个，每个操作又不能拆成多段交给不同核心。因此，除了总工作量，还可以利用操作不可分割这一条件。

将\(S\)中\(n\)个操作的时长降序排列为\(d_{(1)}\ge\cdots\ge d_{(n)}\)。定义

\[A_K(S)=\max_{\substack{r\ge1\\m=(r-1)K+1\le n}}\sum_{j=m-r+1}^{m}d_{(j)}.\]

于是

\[M\ge a+b+\max\left\{\left\lceil\frac{\sum_{v\in S}d_v}{K}\right\rceil,A_K(S)\right\}.\]

\begin{proof} 令\(x=M-a-b\)，\(K\)个同类Pipe在该区间内最多提供\(Kx\)个计算周期，得到第一项。取前\(m=(r-1)K+1\)个操作，至少一个核心要执行其中\(r\)个；即使选其中最短的\(r\)个，也至少需要\(\sum_{j=m-r+1}^{m}d_{(j)}\)周期。对允许的\(r\)取最大得到第二项。两项都限制同一个区间的长度，故取最大值，不相加。最后只加一次前后所需的\(a+b\)。\end{proof}

例如，两个核心执行三个不可拆分操作，时长分别为10、8、7 cycles。总工作量给出13 cycles的下界，上式利用不可分割性给出15 cycles。达到15仍可能受其他依赖和资源限制，这个例子只说明下界为何能加强。

图上递推\(h_v,t_v\)需要\(O(n+m)\)时间；固定\(a,b\)后，检查集合成员需\(O(n)\)，对选中操作排序需\(O(|S|\log|S|)\)。枚举所有\(a,b\)还要另计成本。若\(S\)为空，不能仅凭\(a+b\)得到有效下界。

\subsection{计算工作量与不同核心之间的等待约束}

考虑一串前后有依赖的特定计算操作。要求每个所选操作与其间全部相关计算都存在有向路径上的先后关系，从而将其他计算放入前后两个所选操作之间、互不重叠的时间区间。只有满足这一条件，后文才能把各区间的必要时间相加；任取几个汇合节点并不足够。

对已证明必须经过的分配到不同核心的计算操作之间的依赖，等待常数记为\(\tau_q\)，问题一取1000，问题二、三取500 cycles。某个区间内选取B.2的非空使用同一种Pipe的操作集合，前后必要计算时长为\(a,b\)。令\(\ell,r\in\{0,1\}\)分别表示区间前后是否有所选约束操作。记\((x)_+=\max(x,0)\)，扣除前后计算已经覆盖的时间后，仍需考虑的等待为

\[c(a,b)=\ell(\tau_q-a)_++r(\tau_q-b)_+.\]

设整个时间区间长度为\(x\)，可安排集合中操作的时间为\(z=x-a-b\)，总计算工作量为\(W(S)=\sum_{v\in S}d_v\)。若前后两个所选操作在同一核心，该核心最多可使用\(z\)，其他核心各至多可使用\((z-c)_+\)，故

\[W(S)\le z+(K-1)(z-c)_+.\]

若前后操作位于不同核心，必有\(K\ge2\)。令\(A=(\tau_q-a)_+\)、\(B=(\tau_q-b)_+\)，可供同类计算使用的总周期数至多为

\[(z-A)_++(z-B)_++(K-2)(z-A-B)_+.\]

逐段比较\(z\)与\(A,B,A+B\)可得

\[(z-A)_++(z-B)_+\le z+(z-A-B)_+,\]

所以仍可采用前式限制工作量。对整数周期，满足\(W\le z+(K-1)(z-c)_+\)所需的最小\(z\)为

\[\phi_{K,c}(W)=\min\left\{W,\left\lceil\frac{W+(K-1)c}{K}\right\rceil\right\}.\]

当\(W\le c\)时，\(z=W\)满足第一段；当\(W>c\)时，解第二段线性不等式并向上取整。结合操作不可分割的限制，得到

\[x\ge a+b+\max\{\phi_{K,c(a,b)}(W(S)),A_K(S)\}.\]

满足开头条件时，不同区间及分隔计算互不重叠，可以把各自已证下界相加。如果只是从几个相互重叠的窗口得到下界，就只能取最大值，不能重复累计。问题一的100 cycles在同一个核心中切换Task的等待也不能自动加入，因为方案未必在所选计算处划分Task。

由上述推导可见，与前后计算重叠的等待时间已被扣除，各时间段只计算一次。本节是补充理论推导，正文下界统计使用的是第7章所列已有版本的结果。

\subsection{状态信息充分性的判定条件}

本节讨论一个尚未由本文主算法一般实现的问题：能否只保存方案的一部分信息，就正确判断后续修改及其结果。两份方案当前Makespan相同，并不足以说明这一点。例如，它们可能采用不同的核心顺序，同一次移动操作在一份方案中合法，在另一份中却造成循环依赖。

先固定允许的修改集合\(\mathcal A\)，例如在已经声明的结构范围内移动、合并或拆分子图；每种修改的参数含义也固定。用\(D_a\)表示允许执行修改\(a\)的方案集合，用\(\Phi(P)\)表示需要保留的观测结果。若只研究成功执行，可取观测为Makespan；若还包含失败，则须把成功或失败的状态纳入观测，不能把未定义的Makespan填成一个周期数。

令\(\mathcal L(P)\)为从方案\(P\)出发能够连续执行的所有有限修改序列，包含不作修改的空序列。两份方案在本节所关心的意义下等价，要求

\[\mathcal L(P)=\mathcal L(Q),\qquad \Phi(w(P))=\Phi(w(Q)),\quad\forall w\in\mathcal L(P).\]

其中\(w(P)\)表示按序列\(w\)依次修改\(P\)。第一个条件保证可做的修改完全一致，第二个保证修改后的观测一致；只比较双方碰巧都能做的修改会漏掉第一类差别。

\textbf{命题B.1。} 设\(S(P)\)为准备保存的信息。对每种修改\(a\)，若存在判定函数\(d_a\)、更新函数\(\bar a\)和观测恢复函数\(o\)，使声明范围内的全部方案满足

\[P\in D_a\iff d_a(S(P)),\qquad S(a(P))=\bar a(S(P)),\qquad\Phi(P)=o(S(P)),\]

且更新等式在\(a\)可执行时成立，则\(S(P)=S(Q)\)蕴含上述等价关系。

\begin{proof} 不作修改时，由观测恢复等式知两份观测相同。若执行相同的前\(j\)步后保留的信息仍相同，则\(d_a\)保证下一步在两份方案中同时允许或同时不允许；允许时，\(\bar a\)使更新后的信息仍相同，\(o\)再恢复相同观测。对序列长度归纳，即得全部可执行序列与后续观测一致。\end{proof}

命题给出充分条件，没有证明保存的信息已经最少。要说明最少，还需证明不同保存状态必能被某个后续修改或观测区分。把整份方案编码成一个整数没有减少信息量；重复使用相同编译结果、保存程序运行现场、删除完全相同的候选，则分别解决较小的问题，不能自动获得本命题关于所有后续修改的保证。

\Needspace{14\baselineskip}
\section{固定实现与复现信息}
三题固定求解器的提交分别为下列SHA。逐输入的图、方案与结果来源及完整SHA-256见配套CSV与来源收据。固定版本结果不能用不同算法的历史赢家拼接。
\begin{itemize}
\item P1：\texttt{\seqsplit{834d8c957538ee069c66aadac9509552a4cc69d7}}；求解过程中的评分至多10次E1。
\item P2：\texttt{\seqsplit{c66559a6f8a31ef7b4720e1f7c3c28d61f8dff3f}}；求解过程中的评分至多3次E2。
\item P3：\texttt{\seqsplit{311322b996c0948e8a6a9c7ec6ddfe6ae41fbee1}}；求解过程中的评分至多3次E0。
\end{itemize}
官方代码与配置身份分别为\texttt{\seqsplit{de11a83db8d7c47ed328b15a7df71d613a833b16cd23ee9fe877999578a1ace0}}和\texttt{\seqsplit{dcd10de54b23f8366428fb24e828812b1da9549e6eae4a3c3f38604fe5ae77b9}}。来源收据记录822个固定Git对象；本稿的数据导出仅读取已有原件，没有新增求解器或评价器调用。

\end{document}
